\documentclass[11pt,a4paper]{report}
% Compilation : tectonic Cours_IA_QUBO.tex
% Figures à données : python3 tools/make_figures.py (lit Code/resultats/ia_qubo_*.json,
% produits par les scripts de Code/ia_qubo/, voir Code/ia_qubo/README.md)
% Préambule commun (repris du cours d optimisation bayésienne)
\usepackage{iftex}
\ifPDFTeX
  \usepackage[utf8]{inputenc}
  \usepackage[T1]{fontenc}
  \usepackage{lmodern}
\else
  \usepackage{fontspec}
\fi
\usepackage[french]{babel}
\usepackage[a4paper,margin=2.3cm]{geometry}
\usepackage{amsmath,amssymb,mathtools}
\usepackage{xcolor}
\usepackage{tikz}
\usetikzlibrary{arrows.meta,calc,positioning,shapes.geometric,decorations.pathreplacing,decorations.markings,patterns,babel}
\usepackage{pgfplots}
\pgfplotsset{compat=1.18}
\usepgfplotslibrary{groupplots}
\usepackage{tabularx,booktabs,array}
\usepackage{enumitem}
\setlist{nosep,leftmargin=*}
\usepackage{listings}
\usepackage[font=small,labelfont=bf]{caption}
\usepackage[raggedright]{titlesec}
\usepackage[most]{tcolorbox}
\usepackage[colorlinks=true,linkcolor=blue!50!black,urlcolor=blue!50!black,citecolor=blue!50!black]{hyperref}

\definecolor{cblue}{HTML}{2A78D6}
\definecolor{corange}{HTML}{EB6834}
\definecolor{cgreen}{HTML}{1BAF7A}
\definecolor{cink}{HTML}{1F2328}
\definecolor{cmuted}{HTML}{6B6F76}
\definecolor{cteal}{HTML}{0E8A8A}

% ---------------------------------------------------------------- notations
\newcommand{\dd}{\mathrm{d}}
\newcommand{\Ue}{U_e}
\renewcommand{\Re}{\mathit{Re}}
\newcommand{\vect}[1]{\mathbf{#1}}
\newcommand{\code}[1]{\texttt{#1}}

% ---------------------------------------------------------------- encadrés
\tcbset{coursbox/.style={enhanced,breakable,boxrule=0pt,leftrule=2.5pt,arc=0pt,
  left=6pt,right=6pt,top=4pt,bottom=4pt,fonttitle=\bfseries\small,coltitle=cink,
  attach title to upper={\ }}}
\newtcolorbox{definition}[1][]{coursbox,colback=cblue!5,colframe=cblue,title={Définition#1.}}
\newtcolorbox{resultat}[1][]{coursbox,colback=cgreen!6,colframe=cgreen,title={À retenir#1.}}
\newtcolorbox{projet}[1][]{coursbox,colback=corange!6,colframe=corange,title={Dans notre projet#1.}}
\newtcolorbox{piege}[1][]{coursbox,colback=cmuted!7,colframe=cmuted,title={Piège#1.}}
\newtcolorbox{codebox}[1][]{coursbox,colback=cmuted!3,colframe=cblue!60,title={Dans le code#1.}}

% ---------------------------------------------------------------- exercices
\newcounter{exo}[chapter]
\renewcommand{\theexo}{\thechapter.\arabic{exo}}
\newenvironment{exercice}[1][]{\refstepcounter{exo}\par\medskip\noindent
  \textbf{\color{cblue}Exercice~\theexo}\if\relax\detokenize{#1}\relax\else\ (#1)\fi.\ }{\par}
\newcommand{\exercices}{\section*{Exercices}\addcontentsline{toc}{section}{Exercices}
  \noindent\textit{Corrigés en annexe~\ref{annexe:corriges}. Chercher d'abord sans les regarder.}}

\tikzset{
  flow/.style={draw=cink,rounded corners=2pt,fill=white,align=center,font=\small,inner sep=4pt},
  arr/.style={-{Stealth[length=2.2mm]},cink,thick},
}

\lstset{basicstyle=\ttfamily\footnotesize,breaklines=true,columns=fullflexible,frame=single,rulecolor=\color{cmuted!50},
  backgroundcolor=\color{cmuted!6},keywordstyle=\color{cblue}\bfseries,commentstyle=\color{cmuted}\itshape,
  stringstyle=\color{corange},showstringspaces=false,language=Python,
  literate={é}{\'e}1{è}{\`e}1{à}{\`a}1{ê}{\^e}1{ç}{\c c}1{ô}{\^o}1{î}{\^i}1{ù}{\`u}1{É}{\'E}1{→}{$\to$}1}

\setcounter{secnumdepth}{2}
\setcounter{tocdepth}{1}
\setlength{\parindent}{0pt}
\setlength{\emergencystretch}{3em}
\setlength{\parskip}{0.45em}


\definecolor{cviolet}{HTML}{7A4FB3}
\newtcolorbox{pistes}[1][]{coursbox,colback=cviolet!6,colframe=cviolet,title={Pistes d'amélioration#1.}}
\newcommand{\dB}{\,\mathrm{dB}}
\newcommand{\R}{\mathbb{R}}
\newcommand{\Id}{\mathbf{I}}


% ============================================================================
% Itérations du projet HackathonQuantique : couleurs et marqueurs communs.
% Inclus par tous les cours et par la fiche projet, pour qu'une itération ait
% la même couleur partout. Pour ajouter une itération : une couleur, un nom,
% une date, une ligne dans \legendeiterations.
% ============================================================================
\definecolor{itV0}{HTML}{8D6E63}   % brun   : V0, modèle de la note d'intention
\definecolor{itV1}{HTML}{C2185B}   % framboise : V1, chaîne complète conforme à la note
\definecolor{itV2}{HTML}{00838F}   % cyan foncé : V2, problème reformulé
\definecolor{itV3}{HTML}{E09A00}   % ambre  : V3, concurrent équitable et problème plus grand
\definecolor{itV4}{HTML}{5E35B1}   % indigo : V4, réservé

\makeatletter
\newcommand{\itername}[1]{\@nameuse{itname@#1}}
\newcommand{\iterdate}[1]{\@nameuse{itdate@#1}}
\newcommand{\defiter}[3]{\@namedef{itname@#1}{#2}\@namedef{itdate@#1}{#3}}
\makeatother
\defiter{V0}{Modèle de la note d'intention}{24--30 oct. 2026}
\defiter{V1}{Chaîne complète conforme à la note}{4--5 nov. 2026}
\defiter{V2}{Problème reformulé}{5--6 nov. 2026}
\defiter{V3}{Concurrent équitable, problème plus grand}{à partir du 6 nov. 2026}

% Pastille : \iterbadge{V1}
\newcommand{\iterbadge}[1]{\tikz[baseline=(b.base)]\node[fill=it#1,text=white,rounded corners=2pt,
  inner xsep=3pt,inner ysep=1.5pt,font=\sffamily\bfseries\scriptsize] (b) {#1};}

% Bandeau de début de chapitre : \iterbanner{V0}
\newcommand{\iterbanner}[1]{\par\noindent\begin{tikzpicture}
  \node[fill=it#1!12,draw=it#1,line width=0.6pt,rounded corners=2pt,text width=\dimexpr\linewidth-12pt\relax,
    inner sep=4pt,anchor=west,font=\small] at (0,0)
    {\iterbadge{#1}\ \textbf{\color{it#1}Itération #1} --- \itername{#1} \hfill{\color{it#1}\iterdate{#1}}};
  \end{tikzpicture}\par\medskip}

% Encadré d'évolution : \begin{evolution}{V1}{titre} ... \end{evolution}
\newtcolorbox{evolution}[2]{enhanced,breakable,boxrule=0pt,leftrule=3pt,arc=0pt,
  left=6pt,right=6pt,top=4pt,bottom=4pt,colback=it#1!7,colframe=it#1,
  fonttitle=\bfseries\small,coltitle=it#1,attach title to upper={\ },
  title={\iterbadge{#1}\ Évolution #1 : #2.}}

% Légende des itérations, pour les avant-propos
\newcommand{\legendeiterations}{%
\begin{center}\small
\begin{tabular}{lll}
\iterbadge{V0} & \itername{V0} & \iterdate{V0}\\
\iterbadge{V1} & \itername{V1} & \iterdate{V1}\\
\iterbadge{V2} & \itername{V2} & \iterdate{V2}\\
\iterbadge{V3} & \itername{V3} & \iterdate{V3}\\
\end{tabular}
\end{center}}

% Encadré « pistes d'amélioration » : défini ici s'il manque (cours V0 sans cet encadré)
\providecolor{cviolet}{HTML}{7A4FB3}
\ifcsname pistes\endcsname\else
\newtcolorbox{pistes}[1][]{enhanced,breakable,boxrule=0pt,leftrule=2.5pt,arc=0pt,
  left=6pt,right=6pt,top=4pt,bottom=4pt,fonttitle=\bfseries\small,coltitle=cink,attach title to upper={\ },
  colback=cviolet!6,colframe=cviolet,title={Pistes d'amélioration#1.}}
\fi


\title{\textbf{IA et QUBO}\\[0.4em]\large Apprendre le modèle de la boucle quantique :\\ machine de factorisation, modèle bayésien sur les bits, réseau de neurones,\\ et passage à un problème boîte noire plus grand (lot L6 de la V3)}
\author{Notes de travail, projet \emph{HackathonQuantique}}
\date{Version du 27 novembre 2026}

\begin{document}
\pagenumbering{Alph}
\maketitle
\pagenumbering{roman}

\chapter*{Avant-propos}
\addcontentsline{toc}{chapter}{Avant-propos}
\paragraph{Objet.} La boucle QUBO du projet ajuste un polynôme de degré 2 en variables physiques, puis le minimise sous forme de QUBO. Ce cours remplace ce modèle par des modèles \emph{appris} : une machine de factorisation (méthode FMQA), un modèle quadratique bayésien sur les bits (méthode BOCS) et, pour mesurer le prix de la forme QUBO, un réseau de neurones. Il les compare sur le problème V2, puis sur un problème mis à l'échelle (forme CST à 10 paramètres, 30 bits, $10^9$ profils, solveurs en boîte noire). Il fait enfin le point sur ce que l'intelligence artificielle peut apporter d'autre au projet.

\paragraph{Ce qu'on suppose connu.} Le cours de calcul quantique (QUBO, chapitre~8 ; recuit, chapitre~9 ; pourquoi QAOA n'apporte rien, chapitre~15) et le cours d'optimisation bayésienne (lot L1). Les machines de factorisation, la régression bayésienne et les réseaux de neurones sont posés depuis le début.

\paragraph{Fil conducteur.} La question et le plan d'expérience (ch.~\ref{ch:question}) $\to$ les trois modèles appris (ch.~\ref{ch:fm} à \ref{ch:mlp}) $\to$ leur précision selon le nombre de données (ch.~\ref{ch:precision}) $\to$ le traitement des contraintes (ch.~\ref{ch:contraintes}) $\to$ les résultats sur le problème V2 (ch.~\ref{ch:resultats-v2}) $\to$ le problème mis à l'échelle et ses résultats (ch.~\ref{ch:echelle} et \ref{ch:resultats-echelle}) $\to$ ce que l'IA peut apporter d'autre (ch.~\ref{ch:ia-autres}) $\to$ bilan (ch.~\ref{ch:bilan}).

\paragraph{Code et chiffres.} Tout le code est dans \code{Code/ia\_qubo/}, à part de l'approche actuelle, qu'il ne modifie pas. Les figures et les chiffres sont produits par \code{tools/make\_figures.py} à partir des fichiers \code{Code/resultats/ia\_qubo\_*.json}. Chaque chapitre se termine par des exercices (corrigés en annexe~\ref{annexe:corriges}) et des pistes d'amélioration.

\paragraph{Itérations du projet.} Tout ce cours appartient à l'itération V3 (lot L6).
\legendeiterations

\tableofcontents
\cleardoublepage
\pagenumbering{arabic}

\chapter{La question : apprendre le QUBO, pour des problèmes boîte noire et plus grands}\label{ch:question}
\iterbanner{V3}

\section{Ce que le projet a établi}
La boucle du projet a trois étapes : ajuster un \emph{modèle} sur les profils déjà simulés, choisir le prochain profil en minimisant ce modèle (écrit comme un QUBO), le simuler. Trois résultats sont acquis :
\begin{itemize}
\item la boucle retrouve tout le front traînée/SER en 75 (S1) et 142,5 (S2) simulations, quand NSGA-II n'y parvient presque jamais (V2) ;
\item un modèle plus souple, le krigeage, ne fait pas mieux que le modèle quadratique, et c'est le traitement de la contrainte de manœuvre qui décide du succès (lot L1, cours \emph{Optimisation bayésienne}) ;
\item à 14 qubits, le solveur du QUBO n'est pas le goulot : l'énumération est imbattable, et les pénalités écrasent l'objectif au point d'aveugler QAOA et le recuit (cours de calcul quantique, chapitres 9, 10 et 15).
\end{itemize}

\section{La nouvelle question}
Le modèle de la boucle est un polynôme de degré 2 en variables \emph{physiques} normalisées. C'est un choix fort : il suppose que chaque grandeur varie comme une parabole de chaque paramètre. Il marche sur un NACA à 4 paramètres ; marchera-t-il sur une forme décrite par des dizaines de paramètres, évaluée par des codes dont on ne connaît rien (boîtes noires) ?

L'idée de ce cours est d'utiliser des méthodes d'apprentissage pour \textbf{apprendre directement le QUBO} : un modèle plus souple, ajusté sur les données, mais qui reste quadratique en bits pour pouvoir être confié à un solveur quantique. On en tire trois questions :
\begin{description}
\item[Q1. Précision.] Un modèle appris sur les bits est-il plus précis que le modèle en $z$ ? À partir de combien de données ? Et le gain de précision se traduit-il par moins de simulations ?
\item[Q2. Contraintes.] Peut-on se passer de modèle de contrainte, en laissant le modèle appris \og voir\fg{} les profils infaisables comme mauvais (boîte noire complète) ?
\item[Q3. Passage à l'échelle.] Sur un problème plus grand (30 bits, $10^9$ profils, ni grille ni QUBO énumérables), quelle méthode garde l'avantage ?
\end{description}

\section{Où se situe la recherche}
L'idée n'est pas neuve, et c'est une bonne nouvelle : notre boucle appartient à une famille de méthodes connue, à laquelle on peut se comparer.
\begin{itemize}
\item \textbf{FMQA} (Kitai et al., \emph{Physical Review Research}, 2020) : une machine de factorisation apprend la fonction boîte noire directement sur des bits ; c'est un QUBO par construction ; on le résout par recuit quantique, on simule, on réapprend. Appliquée d'abord à des métamatériaux, étendue depuis aux variables entières (2022), continues (2024), et aux circuits de type QAOA. Un paquet Python existe (\code{fmqa}).
\item \textbf{BOCS} (Baptista et Poloczek, ICML 2018) : modèle quadratique complet sur les bits, régression bayésienne, choix par tirage de Thompson ; le QUBO d'acquisition est résolu par recuit simulé ou programmation semi-définie.
\item \textbf{Espace latent binaire} (Wilson et al., \emph{Applied Physics Reviews}, 2021) : un autoencodeur variationnel binaire apprend un code qui ne décrit que des conceptions valides ; FMQA travaille ensuite dans ce code, et les contraintes n'ont plus besoin de pénalités.
\item \textbf{Applications aéronautiques}, encore rares : une thèse de TU Delft (2025) applique FMQA à des profils couplés à un solveur Lattice Boltzmann ; un article de Keio (2026) l'applique à l'aile d'un avion à propulsion humaine (32 variables) ; un article de Pékin (2026) optimise des profils NACA 4 chiffres avec une machine d'Ising cohérente et des surfaces de réponse.
\end{itemize}
Dans ce vocabulaire, notre boucle est un FMQA dont le modèle n'est pas une machine de factorisation mais un polynôme \emph{structuré} : quadratique dans les variables physiques, ce qui donne un QUBO à couplages contraints ($A^\top MA$).

\section{Le plan d'expérience}
On garde la boucle et l'on change \emph{un seul} élément à la fois : le modèle, puis le traitement de la contrainte.
\begin{center}
\renewcommand{\arraystretch}{1.25}
\small
\begin{tabularx}{\linewidth}{>{\raggedright\arraybackslash}p{3.2cm}>{\raggedright\arraybackslash}X c c}
\toprule
Modèle & En deux mots & QUBO ? & Chapitre \\
\midrule
\code{quadz} & l'approche actuelle : quadratique en $z$ & oui & cours quantique, ch.~8 \\
\code{fm} & machine de factorisation sur les bits (FMQA) & oui & \ref{ch:fm} \\
\code{bocs} & quadratique complet sur les bits, bayésien, tirage de Thompson & oui & \ref{ch:bocs} \\
\code{mlp} & réseau de neurones, recherche directe & non & \ref{ch:mlp} \\
\midrule
Contrainte & & & \\
\midrule
\code{modele} & $C_{L,\max}$ modélisé (affine), variables d'écart, comme en V2 & & \ref{ch:contraintes} \\
\code{bn} (boîte noire) & pas de modèle de contrainte : les infaisables entrent dans les données avec une mauvaise valeur & & \ref{ch:contraintes} \\
\bottomrule
\end{tabularx}
\end{center}
Sur le problème V2, les $4\times2=8$ combinaisons sont comparées avec le protocole du lot L1 : mêmes 10 profils initiaux par graine, 5 poids, un profil par poids et par itération, budget de 250 simulations, 20 graines, nombre de simulations pour retrouver tout le front exact (chapitre~\ref{ch:resultats-v2}). Sur le problème mis à l'échelle, le front exact n'est plus connu : on suit l'hypervolume (chapitres~\ref{ch:echelle} et~\ref{ch:resultats-echelle}).

\section{Le code}
Tout est dans un dossier à part, \code{Code/ia\_qubo/}, qui ne modifie aucun fichier existant : il importe les briques de la V2 (\code{CountedProblem}, \code{z\_of}, les pénalités de \code{qubo\_loop.py}).
\begin{center}
\renewcommand{\arraystretch}{1.2}
\small
\begin{tabularx}{\linewidth}{>{\raggedright\arraybackslash}p{4.6cm}X}
\toprule
Fichier & Contenu \\
\midrule
\code{modeles.py} & les quatre modèles, l'encodage indices $\leftrightarrow$ bits \\
\code{boucle.py} & la boucle générique \code{boucle\_ia}, le recuit vectorisé, la recherche directe \\
\code{experiences\_v2\_ia.py} & la comparaison sur le problème V2 $\to$ \code{resultats/ia\_qubo\_V2.json} \\
\code{precision\_modeles.py} & la précision des modèles selon le nombre de données $\to$ \code{ia\_qubo\_precision.json} \\
\code{probleme\_cst.py} & le problème mis à l'échelle : forme CST, XFOIL et SER en boîte noire, cache \\
\code{experiences\_echelle.py} & la comparaison à l'échelle $\to$ \code{resultats/ia\_qubo\_echelle.json} \\
\bottomrule
\end{tabularx}
\end{center}

\begin{codebox}[ : la boucle générique reproduit la V2]
Avec le modèle \code{quadz}, la contrainte \code{modele} et le solveur exact, \code{boucle\_ia} doit refaire exactement la boucle QUBO de la V2. La fonction \code{verifier\_boucle} le contrôle : sur trois graines, les deux boucles simulent les mêmes profils, dans le même ordre. Toute différence observée ensuite vient donc du seul élément changé.
\end{codebox}

\exercices
\begin{exercice}\label{ex:classer}
Pour chacun des quatre modèles, dire s'il est un QUBO, combien il a de paramètres sur le problème V2 (10 bits, 4 variables), et s'il « sait » que les niveaux d'une variable sont ordonnés.
\end{exercice}
\begin{exercice}\label{ex:fmqa-nous}
En quoi la boucle de la V2 est-elle un cas particulier de FMQA ? Qu'est-ce qui la distingue d'une machine de factorisation ? (Penser à la forme des couplages entre bits.)
\end{exercice}

\begin{pistes}
\begin{itemize}
\item Comparer aussi au paquet \code{fmqa} de référence, pour vérifier que notre machine de factorisation se comporte comme celle de la littérature.
\end{itemize}
\end{pistes}

\chapter{La machine de factorisation : un modèle appris qui est un QUBO}\label{ch:fm}
\iterbanner{V3}

\section{L'idée}
Dans la boucle actuelle, le modèle est un polynôme de degré 2 en variables \emph{physiques} normalisées $z$, que l'on traduit ensuite en bits (cours de calcul quantique, chapitre~8). La machine de factorisation (Rendle, 2010) fait l'inverse : elle apprend \emph{directement} une fonction des bits, déjà sous la forme d'un QUBO. C'est le cœur de FMQA (Kitai et al., 2020).

\begin{definition}[ (machine de factorisation de rang $k$)]
Pour $n$ bits $b_1,\dots,b_n$ :
\[ \hat y(\mathbf b)=w_0+\sum_{i=1}^nw_i\,b_i+\sum_{i<j}\langle\mathbf v_i,\mathbf v_j\rangle\,b_ib_j,\qquad \mathbf v_i\in\R^k,\quad\langle\mathbf v_i,\mathbf v_j\rangle=\sum_{f=1}^kv_{if}v_{jf}. \]
Les paramètres sont $w_0$, les $n$ poids $w_i$ et les $n$ vecteurs $\mathbf v_i$ de dimension $k$ : $1+n+nk$ nombres.
\end{definition}

\paragraph{C'est un QUBO, par construction.} Il suffit de lire les coefficients :
\[ Q_{ii}=w_i,\qquad Q_{ij}=\langle\mathbf v_i,\mathbf v_j\rangle\ (i<j),\qquad c=w_0. \]
Aucune substitution, aucune approximation : l'énergie du QUBO \emph{est} la prédiction du modèle. Le code le vérifie à $10^{-14}$ près.

\paragraph{Pourquoi factoriser les couplages ?} Un modèle quadratique complet sur les bits a un couplage libre $J_{ij}$ par paire, soit $n(n-1)/2$ coefficients estimés chacun séparément. La machine de factorisation impose $J_{ij}=\langle\mathbf v_i,\mathbf v_j\rangle$, c'est-à-dire que la matrice des couplages est de la forme $VV^\top$ (hors diagonale), de rang au plus $k$. Deux conséquences :
\begin{itemize}
\item \textbf{moins de paramètres} : $nk$ au lieu de $n(n-1)/2$ ;
\item \textbf{des couplages qui partagent l'information} : le vecteur $\mathbf v_i$ est appris à partir de \emph{tous} les couplages du bit $i$. Un couplage $J_{ij}$ rarement « observé » (les bits $i$ et $j$ rarement allumés ensemble dans les données) est quand même estimé, par l'intermédiaire de $\mathbf v_i$ et $\mathbf v_j$.
\end{itemize}
C'est l'idée des systèmes de recommandation, pour lesquels la méthode a été inventée : prédire la note d'un utilisateur pour un film qu'il n'a pas vu, à partir de vecteurs appris sur les films qu'il a vus.

\begin{center}
\renewcommand{\arraystretch}{1.2}
\small
\begin{tabular}{lccc}
\toprule
Nombre de paramètres & V2 : 10 bits & CST : 30 bits & 20 variables : 80 bits \\
\midrule
quadratique en $z$ (approche actuelle) & 15 & 66 & 231 \\
machine de factorisation, $k=4$ & 51 & 151 & 401 \\
quadratique complet sur les bits (chapitre~\ref{ch:bocs}) & 56 & 466 & 3\,241 \\
\bottomrule
\end{tabular}
\end{center}

\section{Un exemple à la main}
Trois bits, rang $k=2$ : $w_0=1$, $\mathbf w=(-1,\,0{,}5,\,-2)$, et
\[ \mathbf v_1=(1,\,0),\qquad\mathbf v_2=(0{,}5,\,0{,}5),\qquad\mathbf v_3=(-1,\,1). \]
Les couplages valent $J_{12}=\langle\mathbf v_1,\mathbf v_2\rangle=0{,}5$, $J_{13}=-1$, $J_{23}=-0{,}5+0{,}5=0$. Le QUBO est donc
\[ Q=\begin{pmatrix}-1&0{,}5&-1\\0&0{,}5&0\\0&0&-2\end{pmatrix},\qquad c=1. \]
Pour la chaîne $101$ : $\hat y=1-1-2+J_{13}=-3$. On l'obtient aussi par $\mathbf b^\top Q\mathbf b+c=Q_{11}+Q_{33}+Q_{13}+c=-1-2-1+1=-3$.

\paragraph{Le calcul rapide de la somme des couplages.} Calculer $\sum_{i<j}\langle\mathbf v_i,\mathbf v_j\rangle b_ib_j$ directement coûte $O(n^2k)$. Comme $b_i^2=b_i$ :
\[ \sum_{i<j}\langle\mathbf v_i,\mathbf v_j\rangle b_ib_j=\frac12\sum_{f=1}^k\Big[\Big(\sum_iv_{if}b_i\Big)^2-\sum_iv_{if}^2\,b_i\Big], \]
en $O(nk)$. Pour $101$ : composante $f=1$, $\sum_iv_{i1}b_i=1-1=0$ et $\sum_iv_{i1}^2b_i=1+1=2$, soit $\frac12(0-2)=-1$ ; composante $f=2$, $\sum v_{i2}b_i=0+1=1$ et $\sum v_{i2}^2b_i=0+1=1$, soit $0$. Total $-1=J_{13}$ : on retrouve le couplage.

\section{L'apprentissage}
On ajuste les paramètres en minimisant l'erreur quadratique moyenne sur les $N$ profils simulés, avec une pénalité $L_2$ (comme la régularisation ridge, cours de calcul quantique, chapitre~8) :
\[ \mathcal L=\frac1N\sum_{n=1}^N\big(\hat y(\mathbf b_n)-y_n\big)^2+\lambda\big(\|\mathbf w\|^2+\|V\|^2\big),\qquad\lambda=10^{-3}. \]
Contrairement au modèle quadratique en $z$, ce problème n'est \textbf{pas} un problème de moindres carrés linéaire : $\hat y$ dépend des $\mathbf v_i$ par des produits $v_{if}v_{jf}$. Il n'y a pas de formule fermée, on descend le gradient. En notant $r_n=\hat y(\mathbf b_n)-y_n$ et $s_{nf}=\sum_iv_{if}b_{ni}$ :
\[ \frac{\partial\mathcal L}{\partial w_i}=\frac2N\sum_nr_nb_{ni}+2\lambda w_i,\qquad \frac{\partial\mathcal L}{\partial v_{if}}=\frac2N\sum_nr_n\,b_{ni}\big(s_{nf}-v_{if}\big)+2\lambda v_{if}. \]
(La dérivée de $\frac12[s_f^2-\sum_iv_{if}^2b_i]$ par rapport à $v_{if}$ vaut $b_is_f-v_{if}b_i$.) On utilise l'algorithme Adam, 1\,500 itérations, sur la sortie centrée réduite.

\paragraph{Plusieurs départs, une moyenne.} Le problème n'étant pas convexe, le résultat dépend du tirage initial des $\mathbf v_i$. Le code ajuste 3 machines avec des graines différentes et prend la \emph{moyenne} de leurs QUBO. La moyenne de QUBO est un QUBO : la propriété est conservée.

\begin{codebox}[ : \code{ia\_qubo/modeles.py}, classe \code{FactorizationMachine}]
\begin{lstlisting}
BV = B @ V                                          # s_nf, (N, k)
pred = w0 + B @ w + 0.5 * ((BV ** 2).sum(1) - (B @ (V ** 2)).sum(1))
r = pred - y
gw = 2 * B.T @ r / N + 2 * l2 * w
gV = 2 * (B.T @ (r[:, None] * BV) - (B.T @ r)[:, None] * V) / N + 2 * l2 * V
...
def to_qubo(self):                                  # Q_ii = w_i, Q_ij = <v_i, v_j>
    G = V @ V.T
    Q[np.diag_indices(n)] = w ; Q[iu, ju] = G[iu, ju]
\end{lstlisting}
\end{codebox}

\section{Ce que la machine de factorisation voit, et ce qu'elle ne voit pas}
\paragraph{Plus souple par variable.} Sur les 3 bits d'une variable à 8 niveaux, un modèle quadratique en bits a 1 constante, 3 poids et 3 couplages : 7 degrés de liberté pour 8 niveaux. Il peut suivre presque n'importe quelle courbe en fonction de l'indice. Le modèle en $z$, lui, impose une parabole (3 coefficients). Si la traînée dépendait de la cambrure de façon non parabolique, seul le modèle en bits pourrait le représenter.

\begin{piege}[ : les bits ne connaissent pas l'ordre des niveaux]
En binaire, $x=3$ s'écrit $110$ et $x=4$ s'écrit $001$ (poids faible à gauche) : les deux niveaux voisins n'ont \emph{aucun} bit allumé en commun. Pour le modèle en bits, 3 et 4 ne sont pas plus proches que 0 et 7. Il ne sait pas que la traînée varie continûment avec la cambrure : il doit l'\emph{apprendre} à partir des données. Le modèle en $z$, au contraire, a la continuité dans sa forme même : $z$ est proportionnel à l'indice. Avec peu de données, cette connaissance a priori vaut beaucoup (chapitre~\ref{ch:precision}).
\end{piege}

\begin{resultat}
\begin{itemize}
\item La machine de factorisation est un modèle quadratique en bits dont les couplages sont $\langle\mathbf v_i,\mathbf v_j\rangle$ : c'est un QUBO, sans traduction.
\item Elle a $1+n+nk$ paramètres, entre le modèle en $z$ (le plus contraint) et le quadratique complet en bits (le plus libre).
\item Elle s'ajuste par descente de gradient (problème non convexe) ; on moyenne plusieurs ajustements.
\item Elle peut représenter des effets non paraboliques, mais elle ignore que les niveaux sont ordonnés : il lui faut plus de données.
\end{itemize}
\end{resultat}

\exercices
\begin{exercice}\label{ex:fm-main}
Avec la machine de l'exemple, calculer $\hat y$ pour les 8 chaînes de 3 bits, par la formule du QUBO. Quelle chaîne est le minimum ? Vérifier une valeur avec la formule rapide en $O(nk)$.
\end{exercice}
\begin{exercice}\label{ex:fm-rang}
(a) Montrer qu'aucune machine de rang $k=1$ ne peut donner $J_{12}=J_{13}=J_{23}=-1$. (Indice : calculer le produit $J_{12}J_{13}J_{23}$.) (b) Trouver trois vecteurs de dimension $k=2$ qui donnent ces trois couplages. (Indice : trois vecteurs de même norme à $120^\circ$ les uns des autres.) (c) Qu'en conclure sur le choix de $k$ ?
\end{exercice}
\begin{exercice}\label{ex:fm-ordre}
Sur une variable à 8 niveaux codée en binaire sur 3 bits, combien de bits diffèrent entre les niveaux 3 et 4 ? Entre 0 et 7 ? Entre 6 et 7 ? Pourquoi un modèle en bits entraîné sans aucun profil au niveau 4 peut-il prédire n'importe quoi à ce niveau, alors que le modèle en $z$ l'interpole ?
\end{exercice}

\begin{pistes}
\begin{itemize}
\item Encoder les niveaux en \og thermomètre\fg{} (unaire : $x=\sum_jb_j$) plutôt qu'en binaire : deux niveaux voisins diffèrent alors d'un seul bit, et le modèle en bits retrouve une notion d'ordre, au prix de 7 bits par variable à 8 niveaux au lieu de 3.
\item Régler le rang $k$ et la pénalité $\lambda$ par validation croisée à chaque itération, plutôt que de les fixer.
\end{itemize}
\end{pistes}

\chapter{Le modèle bayésien sur les bits et le tirage de Thompson}\label{ch:bocs}
\iterbanner{V3}

\section{Le modèle le plus libre qui reste un QUBO}
BOCS (\emph{Bayesian Optimization of Combinatorial Structures}, Baptista et Poloczek, 2018) prend le modèle quadratique \emph{complet} sur les bits : une constante, un poids par bit, un couplage libre par paire.
\[ \hat y(\mathbf b)=\boldsymbol\varphi(\mathbf b)\cdot\boldsymbol\theta,\qquad\boldsymbol\varphi(\mathbf b)=\big(1;\ b_1,\dots,b_n;\ b_1b_2,\ b_1b_3,\dots,b_{n-1}b_n\big), \]
soit $p=1+n+n(n-1)/2$ coefficients : 56 pour nos 10 bits, 466 pour les 30 bits du problème mis à l'échelle. Comme pour la machine de factorisation, les coefficients se lisent directement comme un QUBO : $Q_{ii}=\theta_{b_i}$, $Q_{ij}=\theta_{b_ib_j}$, $c=\theta_1$ (\code{theta\_to\_qubo}).

La différence est dans la façon d'ajuster : au lieu d'un seul jeu de coefficients, on calcule une \emph{loi de probabilité} sur les coefficients, qui dit ce que les données permettent de savoir et ce qu'elles laissent incertain.

\section{La régression bayésienne}
\paragraph{Les hypothèses.} A priori (avant les données), chaque coefficient suit une loi normale centrée : $\theta_k\sim\mathcal N(0,1/\alpha)$, sauf la constante, laissée libre. Les simulations sont le modèle plus un bruit : $y_n=\boldsymbol\varphi(\mathbf b_n)\cdot\boldsymbol\theta+\varepsilon_n$, $\varepsilon_n\sim\mathcal N(0,1/\beta)$. On travaille sur $y$ centré réduit.

\paragraph{La loi a posteriori.} Le produit de deux gaussiennes est une gaussienne : après les données, $\boldsymbol\theta\sim\mathcal N(\mathbf m,S)$ avec
\[ S=\big(A+\beta\,\Phi^\top\Phi\big)^{-1},\qquad\mathbf m=\beta\,S\,\Phi^\top\mathbf y,\qquad A=\mathrm{diag}(\epsilon,\alpha,\dots,\alpha). \]
$\Phi$ est la matrice $N\times p$ des $\boldsymbol\varphi(\mathbf b_n)$, et $\epsilon=10^{-6}$ laisse la constante presque libre.

\paragraph{Le lien avec la régularisation ridge.} En multipliant par $\beta$ : $\mathbf m=(\Phi^\top\Phi+A/\beta)^{-1}\Phi^\top\mathbf y$. C'est exactement la solution ridge, avec $\lambda=\alpha/\beta$, et \emph{sans pénaliser la constante} : la correction que le cours de calcul quantique (chapitre~8, section sur la régularisation) recommandait pour le modèle en $z$. La moyenne a posteriori est le modèle ridge ; la covariance $S$ dit en plus à quel point chaque coefficient est incertain.

\paragraph{Régler $\alpha$ et $\beta$ sur les données.} Plutôt que de les fixer, on les choisit pour que les données soient les plus probables possible sous le modèle (maximum de l'\emph{évidence}, MacKay 1992). On itère trois formules jusqu'à convergence :
\[ \gamma=\sum_k\frac{\beta\lambda_k}{\alpha+\beta\lambda_k},\qquad\alpha=\frac{\gamma}{\|\mathbf m\|^2},\qquad\beta=\frac{N-\gamma}{\|\mathbf y-\Phi\mathbf m\|^2}, \]
où les $\lambda_k$ sont les valeurs propres de $\Phi^\top\Phi$. Le nombre $\gamma$ a un sens simple : c'est le \textbf{nombre effectif de coefficients} que les données déterminent vraiment. Chaque terme vaut presque 1 dans une direction bien informée ($\beta\lambda_k\gg\alpha$), presque 0 dans une direction mal informée. Avec 30 profils et 56 coefficients, $\gamma$ ne peut pas dépasser 30 : les 26 autres directions restent gouvernées par l'a priori.

\section{Le tirage de Thompson : choisir avec l'incertitude}
Comment choisir le prochain profil avec une loi sur les modèles plutôt qu'un modèle ? L'optimisation bayésienne du lot L1 utilisait l'amélioration espérée, qui n'est pas quadratique : on ne peut pas l'écrire comme un QUBO. BOCS utilise une règle qui, elle, le permet :
\begin{enumerate}
\item tirer \emph{un} jeu de coefficients au hasard dans la loi a posteriori : $\tilde{\boldsymbol\theta}\sim\mathcal N(\mathbf m,S)$ ;
\item ce tirage est un modèle quadratique en bits, donc un QUBO : le minimiser ;
\item simuler le profil obtenu.
\end{enumerate}
C'est l'échantillonnage de Thompson (1933). Il équilibre de lui-même exploration et exploitation : un coefficient bien connu est tiré près de sa moyenne, un coefficient incertain peut être tiré très favorable, ce qui envoie la boucle explorer la région qu'il gouverne. Quand les données s'accumulent, $S$ rétrécit et les tirages se resserrent autour de $\mathbf m$.

\paragraph{Exemple chiffré.} Supposons qu'un couplage ait pour moyenne a posteriori $m=-0{,}2$ et pour écart-type $0{,}5$ (incertain). Il est tiré négatif, donc favorable à l'allumage simultané des deux bits, avec probabilité $\Phi(0{,}2/0{,}5)=\Phi(0{,}4)\approx0{,}66$, et positif une fois sur trois. Avec 20 fois plus de données, l'écart-type tombe vers $0{,}5/\sqrt{20}\approx0{,}11$ et la probabilité monte à $\Phi(1{,}8)\approx0{,}96$ : le doute est levé.

\paragraph{Le tirage en pratique.} Avec la décomposition de Cholesky $S=LL^\top$ et $\boldsymbol\varepsilon\sim\mathcal N(0,\Id)$, le vecteur $\tilde{\boldsymbol\theta}=\mathbf m+L\boldsymbol\varepsilon$ suit la bonne loi. Dans la boucle, chaque poids $w$ tire ses propres modèles de traînée et de SER : les cinq QUBO d'une itération sont cinq tirages différents.

\begin{codebox}[ : \code{ia\_qubo/modeles.py}, classe \code{BayesQuadBits}]
\begin{lstlisting}
for _ in range(self.n_iter):                     # maximum de l'evidence (MacKay)
    S = np.linalg.inv(A + beta * G)               # G = Phi^T Phi
    mN = beta * S @ Phi.T @ t
    gamma = np.sum(beta * lam / (alpha + beta * lam))
    alpha = gamma / (mN[1:] @ mN[1:])            # la constante n'entre pas dans la pénalité
    beta = (N - gamma) / np.sum((t - Phi @ mN) ** 2)
...
def sample_qubo(self):                            # tirage de Thompson -> QUBO
    L = np.linalg.cholesky(self.S + 1e-10 * np.eye(len(self.m)))
    return self._to_qubo(self.m + L @ self.rng.standard_normal(len(self.m)))
\end{lstlisting}
\end{codebox}

\begin{piege}[ : le prix de la liberté]
Le modèle complet a $\Theta(n^2)$ coefficients. À 80 bits, il en a 3\,241 : il faudrait des milliers de simulations pour les déterminer, alors que l'on en a quelques centaines. L'a priori gaussien ne crée pas d'information : il ramène vers 0 tout ce que les données ne déterminent pas. BOCS d'origine utilise un a priori \og en fer à cheval\fg{} (\emph{horseshoe}), qui suppose que la plupart des couplages sont nuls ; c'est une hypothèse plus forte, utile quand elle est vraie.
\end{piege}

\begin{resultat}
\begin{itemize}
\item BOCS ajuste le modèle quadratique complet en bits par régression bayésienne : la moyenne a posteriori est la solution ridge sans pénaliser la constante, et la covariance mesure l'incertitude de chaque coefficient.
\item Le tirage de Thompson remplace l'amélioration espérée : un modèle tiré au hasard dans la loi a posteriori est un QUBO, et le minimiser explore là où le modèle doute.
\item C'est le plus libre des modèles quadratiques en bits, donc le plus gourmand en données.
\end{itemize}
\end{resultat}

\exercices
\begin{exercice}\label{ex:bocs-ridge}
Montrer que $\mathbf m=\beta S\Phi^\top\mathbf y$ est la solution de $\min_{\boldsymbol\theta}\ \beta\|\Phi\boldsymbol\theta-\mathbf y\|^2+\boldsymbol\theta^\top A\boldsymbol\theta$. En déduire la valeur du paramètre ridge $\lambda$ équivalent.
\end{exercice}
\begin{exercice}\label{ex:thompson}
Un poids de bit a pour loi a posteriori $\mathcal N(0{,}3;\ 0{,}2^2)$. (a) Probabilité qu'un tirage de Thompson le donne négatif (bit favorable à allumer) ? (b) Même question pour $\mathcal N(0{,}3;\ 1^2)$. (c) Expliquer pourquoi le tirage de Thompson explore davantage en début de boucle.
\end{exercice}
\begin{exercice}\label{ex:gamma}
Avec $N=30$ profils et le modèle complet à 56 coefficients, pourquoi a-t-on toujours $\gamma\le30$ ? Que vaut à peu près $\gamma$ pour un modèle en $z$ à 15 coefficients, si $\Phi^\top\Phi$ a toutes ses valeurs propres grandes devant $\alpha/\beta$ ?
\end{exercice}

\begin{pistes}
\begin{itemize}
\item Utiliser l'a priori en fer à cheval de BOCS d'origine, qui favorise les modèles creux, et comparer.
\item Appliquer le tirage de Thompson au modèle en $z$ (régression bayésienne sur ses 15 coefficients) : on garderait la continuité des niveaux, avec en plus l'exploration.
\end{itemize}
\end{pistes}

\chapter{Le réseau de neurones : précis, mais pas un QUBO}\label{ch:mlp}
\iterbanner{V3}

\section{Pourquoi l'étudier}
Les trois modèles précédents sont quadratiques en bits, parce que le solveur quantique l'exige. Un réseau de neurones ne l'est pas. Le mettre dans la même boucle répond à une question simple : \textbf{combien coûte la contrainte « le modèle doit être un QUBO »} ? Si un modèle libre fait nettement mieux, c'est autant de précision que la formulation quantique oblige à abandonner.

\section{Le modèle}
\begin{definition}[ (perceptron à deux couches cachées)]
Pour une entrée $\mathbf z\in\R^d$ (les variables normalisées, comme le modèle de la V2) :
\[ \mathbf a_1=\tanh(W_1^\top\mathbf z+\mathbf c_1),\qquad\mathbf a_2=\tanh(W_2^\top\mathbf a_1+\mathbf c_2),\qquad\hat y=\mathbf w_3\cdot\mathbf a_2+c_3, \]
avec 32 neurones par couche cachée. La fonction $\tanh$ est appliquée composante par composante.
\end{definition}
Pour $d=4$, cela fait $4\times32+32+32\times32+32+32+1=1\,249$ paramètres, presque cent fois plus que le modèle quadratique en $z$ (15). Pour $d=10$ (problème mis à l'échelle), 1\,441.

\paragraph{L'apprentissage.} Comme pour la machine de factorisation : erreur quadratique moyenne plus pénalité $L_2$ ($\lambda=10^{-3}$), sur la sortie centrée réduite, minimisée par Adam (1\,000 à 1\,500 itérations). Le gradient se calcule couche par couche, de la sortie vers l'entrée (rétropropagation) : avec $\mathbf r=\frac2N(\hat{\mathbf y}-\mathbf y)$,
\[ \frac{\partial\mathcal L}{\partial\mathbf w_3}=A_2^\top\mathbf r,\qquad\boldsymbol\delta_2=(\mathbf r\,\mathbf w_3^\top)\odot(1-A_2^2),\qquad\frac{\partial\mathcal L}{\partial W_2}=A_1^\top\boldsymbol\delta_2,\qquad\boldsymbol\delta_1=(\boldsymbol\delta_2W_2^\top)\odot(1-A_1^2),\ \dots \]
où $1-\tanh^2$ est la dérivée de $\tanh$ et $\odot$ le produit composante par composante.

\paragraph{Un ensemble, pour avoir une incertitude.} Le code entraîne 3 à 5 réseaux à partir de tirages initiaux différents. La prédiction est leur moyenne ; leur écart-type donne une estimation, grossière, de l'incertitude. Là où les données sont nombreuses, les réseaux sont d'accord ; ailleurs, ils divergent.

\section{Pourquoi ce n'est pas un QUBO}
Une fonction de bits s'écrit d'une seule façon comme un polynôme multilinéaire (chaque bit au plus à la puissance 1). Le modèle en $z$ et les modèles en bits n'ont que des termes de degré 0, 1 et 2. Le réseau, lui, compose des $\tanh$ : son développement contient des termes de \emph{tous} les degrés.

\paragraph{Exemple : une seule variable.} Prenons $\hat y=\tanh(3z_m)$ pour la cambrure (8 niveaux, 3 bits). Les 8 valeurs, de $z=-1$ à $z=1$, sont
\[ -0{,}995;\ -0{,}973;\ -0{,}858;\ -0{,}404;\ 0{,}404;\ 0{,}858;\ 0{,}973;\ 0{,}995. \]
Une fonction de 3 bits a 8 coefficients multilinéaires : constante, 3 linéaires, 3 couplages et \emph{un} terme cubique $b_0b_1b_2$. On le trouve en combinant les 8 valeurs avec des signes alternés (exercice~\ref{ex:cubique}) : il ne vaut pas 0 ici. Ce terme n'a pas de place dans un QUBO ; il faudrait un qubit auxiliaire (réduction de Rosenberg) par terme de ce type.

\section{Choisir sans QUBO : la recherche directe}
Sans QUBO, ni l'énumération des chaînes de bits, ni le recuit, ni QAOA ne s'appliquent tels quels. On minimise directement la prédiction sur la grille des profils :
\begin{itemize}
\item \textbf{problème V2} (1\,024 profils) : on évalue le réseau sur toute la grille, ce qui prend quelques millisecondes ;
\item \textbf{problème mis à l'échelle} ($8^{10}\approx10^9$ profils) : on tire 20\,000 profils au hasard, puis, depuis les 20 meilleurs, une \textbf{recherche locale} : on essaie de changer le niveau d'une seule variable (10 variables $\times$ 7 autres niveaux $=70$ voisins) et l'on garde le meilleur voisin, jusqu'à ce qu'aucun ne fasse mieux.
\end{itemize}
Les contraintes sont traitées comme dans le QUBO, sans les défauts de l'encodage : un réseau de plus prédit $C_{L,\max}$, et l'on ajoute $75\,\max(0,\ 1{,}2-\widehat{C_{L,\max}})^2$ (la pénalité idéale, sans les dents de scie des variables d'écart) ; l'épaisseur interdite coûte 3.

\begin{codebox}[ : \code{ia\_qubo/boucle.py}, fonction \code{recherche\_directe}]
\begin{lstlisting}
I = np.column_stack([rng.integers(0, L, n_alea) for L in levels])   # 20 000 profils au hasard
s = score(I)
for x in [I[j] for j in np.argsort(s)[:n_depart]]:                 # 20 departs
    while amelio:                                                   # recherche locale
        voisins = [x avec la variable i au niveau v, pour tous i et v != x[i]]
        j = argmin(score(voisins))
        if score(voisins[j]) < f(x): x = voisins[j]
\end{lstlisting}
\end{codebox}

\begin{projet}
La recherche directe n'a rien de quantique, et rien ne garantit qu'elle trouve le minimum du réseau sur $10^9$ profils. Mais elle est rapide, et l'on n'a pas besoin du minimum exact : on cherche un \emph{bon candidat}, que la simulation départagera. C'est exactement l'argument que l'on oppose à QAOA dans le cours de calcul quantique : la précision du choix compte moins que la qualité du modèle.
\end{projet}

\begin{resultat}
\begin{itemize}
\item Un perceptron à deux couches de 32 neurones a environ 1\,250 paramètres : il représente des effets de tous les ordres, mais demande des données.
\item Ce n'est pas un QUBO : son développement en bits contient des termes de degré 3 et plus.
\item On le minimise par énumération (petite grille) ou par tirage aléatoire et recherche locale (grande grille).
\item Il sert de référence : l'écart entre lui et les modèles quadratiques mesure le prix de la formulation QUBO.
\end{itemize}
\end{resultat}

\exercices
\begin{exercice}\label{ex:mlp-param}
Compter les paramètres d'un perceptron à deux couches de $h$ neurones pour $d$ entrées. Application : $d=4$ et $d=10$, $h=32$. Pour quel $h$ le réseau a-t-il autant de paramètres que le modèle quadratique en $z$ à $d=10$ (66) ?
\end{exercice}
\begin{exercice}\label{ex:cubique}
Pour une fonction $f$ de 3 bits, le coefficient du terme $b_0b_1b_2$ de son polynôme multilinéaire vaut $\sum_{\mathbf b}(-1)^{3-|\mathbf b|}f(\mathbf b)$, où $|\mathbf b|$ est le nombre de bits à 1. (a) Le vérifier sur $f=b_0b_1b_2$. (b) Le calculer pour $\hat y=\tanh(3z_m)$ avec les valeurs du cours (attention : la chaîne $b_0b_1b_2$ correspond à l'indice $x=b_0+2b_1+4b_2$). (c) Que vaut-il pour $\hat y=z_m^2$ ?
\end{exercice}

\begin{pistes}
\begin{itemize}
\item \textbf{Distiller le réseau en QUBO} : ajuster une machine de factorisation sur les prédictions du réseau en des milliers de profils tirés au hasard (des données gratuites), puis résoudre ce QUBO. On garderait la précision du réseau là où une forme quadratique peut la représenter.
\item Remplacer l'ensemble par une vraie méthode d'incertitude (abandon aléatoire, réseaux bayésiens) et l'utiliser pour explorer, comme le tirage de Thompson.
\end{itemize}
\end{pistes}

\chapter{La précision des modèles selon le nombre de données}\label{ch:precision}
\iterbanner{V3}

\section{Ce qu'on mesure}
Avant de mettre les modèles dans la boucle, on mesure leur précision seule, avec le script \code{precision\_modeles.py}. Pour chacune des trois grandeurs ($C_D$, SER du secteur S2, $C_{L,\max}$), pour chaque modèle et pour chaque taille $N$ de jeu d'entraînement de 10 à 500, on tire 10 jeux de $N$ profils au hasard sur la grille V2, on ajuste, et l'on mesure sur tous les autres profils :
\begin{itemize}
\item l'\textbf{erreur relative} : l'erreur quadratique moyenne divisée par l'écart-type de la grandeur. Elle vaut 0 pour un modèle parfait et 1 pour un modèle qui prédirait toujours la moyenne ;
\item la \textbf{corrélation de rang} de Spearman entre prédictions et vraies valeurs : 1 si le modèle classe les profils exactement dans le bon ordre. Pour choisir le prochain profil, c'est le classement qui compte.
\end{itemize}
On rapporte la médiane des 10 tirages. Les modèles de traînée et de SER sont ajustés sur les profils faisables (sans le profil aberrant à plus de 200 counts), celui de $C_{L,\max}$ sur les profils où XFOIL a convergé.

\begin{figure}[htbp]
\centering
% figure générée par tools/make_figures.py
\begin{tikzpicture}
\begin{groupplot}[group style={group size=3 by 2,horizontal sep=1.15cm,vertical sep=2.3cm},width=0.33\linewidth,height=4.6cm,
  xmode=log,xlabel={profils d'entraînement $N$},tick label style={font=\scriptsize},label style={font=\scriptsize},
  title style={font=\small},legend style={font=\scriptsize,draw=none,at={(-0.75,-0.45)},anchor=north,legend columns=4}]
\nextgroupplot[title={traînée $C_D$},ylabel={erreur relative},ymin=0,ymax=1.2]
\addplot[cblue,very thick,mark=*,mark size=1.2pt] coordinates {(10,1.0733) (15,1.3907) (20,0.65927) (30,0.52798) (50,0.41384) (80,0.38052) (120,0.35845) (200,0.35249) (300,0.35353) (500,0.36076)};
\addplot[corange,very thick,mark=square*,mark size=1.2pt] coordinates {(10,0.96921) (15,0.90006) (20,0.87233) (30,0.67583) (50,0.6483) (80,0.4589) (120,0.38695) (200,0.36003) (300,0.3421) (500,0.3447)};
\addplot[cgreen,very thick,mark=triangle*,mark size=1.2pt] coordinates {(10,0.94668) (15,0.89333) (20,0.82213) (30,0.67962) (50,0.60391) (80,0.47934) (120,0.41478) (200,0.36929) (300,0.34324) (500,0.35042)};
\addplot[cviolet,very thick,mark=diamond*,mark size=1.2pt] coordinates {(10,0.76051) (15,0.61865) (20,0.59871) (30,0.47948) (50,0.44549) (80,0.31752) (120,0.25977) (200,0.22991) (300,0.20107) (500,0.21607)};
\nextgroupplot[title={SER},ylabel={erreur relative},ymin=0,ymax=1.2]
\addplot[cblue,very thick,mark=*,mark size=1.2pt] coordinates {(10,0.73271) (15,0.30978) (20,0.29751) (30,0.17479) (50,0.14709) (80,0.13906) (120,0.13262) (200,0.12619) (300,0.12733) (500,0.12679)};
\addplot[corange,very thick,mark=square*,mark size=1.2pt] coordinates {(10,0.80948) (15,0.59477) (20,0.55345) (30,0.57149) (50,0.34562) (80,0.22421) (120,0.18872) (200,0.14723) (300,0.14283) (500,0.13811)};
\addplot[cgreen,very thick,mark=triangle*,mark size=1.2pt] coordinates {(10,0.83039) (15,0.59892) (20,0.57279) (30,0.49819) (50,0.41789) (80,0.23868) (120,0.15831) (200,0.13681) (300,0.13128) (500,0.12472)};
\addplot[cviolet,very thick,mark=diamond*,mark size=1.2pt] coordinates {(10,0.28606) (15,0.23842) (20,0.19862) (30,0.13962) (50,0.10501) (80,0.087025) (120,0.083944) (200,0.077794) (300,0.07862) (500,0.074123)};
\nextgroupplot[title={$C_{L,\max}$},ylabel={erreur relative},ymin=0,ymax=1.2]
\addplot[cblue,very thick,mark=*,mark size=1.2pt] coordinates {(10,1.3935) (15,2.3452) (20,1.0533) (30,0.70457) (50,0.60046) (80,0.56493) (120,0.55718) (200,0.54449) (300,0.53785) (500,0.52145)};
\addplot[corange,very thick,mark=square*,mark size=1.2pt] coordinates {(10,1.1011) (15,1.3323) (20,1.1868) (30,1.1968) (50,1.2171) (80,0.85638) (120,0.65153) (200,0.59295) (300,0.56929) (500,0.52862)};
\addplot[cgreen,very thick,mark=triangle*,mark size=1.2pt] coordinates {(10,0.97587) (15,0.95218) (20,0.84813) (30,0.8249) (50,0.79429) (80,0.7558) (120,0.74138) (200,0.68341) (300,0.582) (500,0.54358)};
\addplot[cviolet,very thick,mark=diamond*,mark size=1.2pt] coordinates {(10,0.78367) (15,0.80826) (20,0.74305) (30,0.72409) (50,0.76588) (80,0.70864) (120,0.64511) (200,0.59028) (300,0.53071) (500,0.49408)};
\nextgroupplot[title={traînée $C_D$},ylabel={corrélation de rang},ymin=0,ymax=1]
\addplot[cblue,very thick,mark=*,mark size=1.2pt] coordinates {(10,0.69981) (15,0.65643) (20,0.84395) (30,0.88531) (50,0.91906) (80,0.93427) (120,0.93668) (200,0.93991) (300,0.935) (500,0.93804)};
\addplot[corange,very thick,mark=square*,mark size=1.2pt] coordinates {(10,0.52042) (15,0.54132) (20,0.62969) (30,0.76582) (50,0.80404) (80,0.88674) (120,0.92772) (200,0.93819) (300,0.94048) (500,0.94519)};
\addplot[cgreen,very thick,mark=triangle*,mark size=1.2pt] coordinates {(10,0.43294) (15,0.54284) (20,0.61188) (30,0.73096) (50,0.7947) (80,0.88265) (120,0.91582) (200,0.93503) (300,0.93969) (500,0.94384)};
\addplot[cviolet,very thick,mark=diamond*,mark size=1.2pt] coordinates {(10,0.7625) (15,0.828) (20,0.83756) (30,0.89475) (50,0.90534) (80,0.95349) (120,0.97226) (200,0.97679) (300,0.98147) (500,0.98346)};
\nextgroupplot[title={SER},ylabel={corrélation de rang},ymin=0,ymax=1]
\addplot[cblue,very thick,mark=*,mark size=1.2pt] coordinates {(10,0.78316) (15,0.97092) (20,0.96878) (30,0.99145) (50,0.99202) (80,0.99383) (120,0.9941) (200,0.9949) (300,0.99425) (500,0.99488)};
\addplot[corange,very thick,mark=square*,mark size=1.2pt] coordinates {(10,0.66949) (15,0.87778) (20,0.86055) (30,0.86811) (50,0.95249) (80,0.97938) (120,0.98618) (200,0.99328) (300,0.99317) (500,0.99389)};
\addplot[cgreen,very thick,mark=triangle*,mark size=1.2pt] coordinates {(10,0.61107) (15,0.85916) (20,0.85871) (30,0.90444) (50,0.94015) (80,0.97539) (120,0.98932) (200,0.99343) (300,0.99377) (500,0.99444)};
\addplot[cviolet,very thick,mark=diamond*,mark size=1.2pt] coordinates {(10,0.96274) (15,0.97938) (20,0.98231) (30,0.99255) (50,0.99506) (80,0.99717) (120,0.99701) (200,0.99765) (300,0.99744) (500,0.99788)};
\nextgroupplot[title={$C_{L,\max}$},ylabel={corrélation de rang},ymin=0,ymax=1]
\addplot[cblue,very thick,mark=*,mark size=1.2pt] coordinates {(10,0.34392) (15,0.57935) (20,0.61601) (30,0.78553) (50,0.87559) (80,0.90069) (120,0.91098) (200,0.91568) (300,0.92013) (500,0.92277)};
\addplot[corange,very thick,mark=square*,mark size=1.2pt] coordinates {(10,0.3353) (15,0.52869) (20,0.52041) (30,0.48987) (50,0.46047) (80,0.68989) (120,0.8492) (200,0.89123) (300,0.91147) (500,0.92586)};
\addplot[cgreen,very thick,mark=triangle*,mark size=1.2pt] coordinates {(10,0.54544) (15,0.64613) (20,0.7337) (30,0.75124) (50,0.77832) (80,0.80481) (120,0.82126) (200,0.85946) (300,0.90198) (500,0.91975)};
\addplot[cviolet,very thick,mark=diamond*,mark size=1.2pt] coordinates {(10,0.75561) (15,0.78924) (20,0.82737) (30,0.83337) (50,0.79037) (80,0.81237) (120,0.89217) (200,0.89316) (300,0.9229) (500,0.93582)};
\legend{quadratique en $z$, machine de factorisation, bayésien sur les bits, réseau de neurones}
\end{groupplot}
\end{tikzpicture}

\caption{Précision des quatre modèles sur la grille V2, en fonction du nombre de profils d'entraînement (échelle logarithmique, médiane de 10 tirages). En haut, l'erreur relative (plus bas est meilleur) ; en bas, la corrélation de rang (plus haut est meilleur).}\label{fig:precision-v2}
\end{figure}

\section{Ce qu'on observe}
\begin{center}
\renewcommand{\arraystretch}{1.2}
\small
\begin{tabular}{l cccc cccc}
\toprule
& \multicolumn{4}{c}{erreur relative} & \multicolumn{4}{c}{corrélation de rang} \\
$C_D$, $N=$ & 15 & 30 & 80 & 300 & 15 & 30 & 80 & 300 \\
\midrule
quadratique en $z$ & 1,39 & 0,53 & 0,38 & 0,35 & 0,66 & 0,89 & 0,93 & 0,93 \\
machine de factorisation & 0,90 & 0,68 & 0,46 & 0,34 & 0,54 & 0,77 & 0,89 & 0,94 \\
bayésien sur les bits & 0,89 & 0,68 & 0,48 & 0,34 & 0,54 & 0,73 & 0,88 & 0,94 \\
réseau de neurones & 0,62 & 0,48 & 0,32 & 0,20 & 0,83 & 0,89 & 0,95 & 0,98 \\
\midrule
$C_{L,\max}$, $N=$ & 15 & 30 & 80 & 300 & 15 & 30 & 80 & 300 \\
\midrule
quadratique en $z$ & 2,35 & 0,70 & 0,56 & 0,54 & 0,58 & 0,79 & 0,90 & 0,92 \\
machine de factorisation & 1,33 & 1,20 & 0,86 & 0,57 & 0,53 & 0,49 & 0,69 & 0,91 \\
réseau de neurones & 0,81 & 0,72 & 0,71 & 0,53 & 0,79 & 0,83 & 0,81 & 0,92 \\
\bottomrule
\end{tabular}
\end{center}

\paragraph{1. Le réseau de neurones est le plus précis, même avec peu de données.} Sur la traînée et la SER, il a la plus petite erreur et le meilleur classement à presque toutes les tailles. Sur la SER, son erreur relative vaut 0,29 dès 10 profils, contre 0,73 pour le modèle en $z$. La régularisation $L_2$ et l'ensemble de réseaux le protègent du surapprentissage, et son entrée $z$ lui donne la continuité des niveaux.

\paragraph{2. Les modèles en bits sont en retard tant que les données sont rares.} La machine de factorisation et le modèle bayésien sur les bits ne rejoignent le modèle en $z$ qu'à partir de 80 à 120 profils. Au-delà, les trois modèles quadratiques se valent (erreur relative 0,34 à 0,36 sur $C_D$ à 300 profils) : ils ont atteint ce qu'une forme quadratique peut représenter. La différence se joue donc entièrement au début de la boucle, là où elle compte.

\paragraph{3. Le modèle en $z$ a un mauvais moment à 15 profils.} Son erreur sur $C_D$ vaut 1,07 à 10 profils, \textbf{1,39 à 15}, puis 0,66 à 20. C'est le régime « juste déterminé » décrit dans le cours de calcul quantique (chapitre~8, régularisation) : avec autant de profils que de coefficients, le modèle interpole aussi le bruit avec de grands coefficients. Les modèles en bits, plus régularisés, n'ont pas ce pic.

\paragraph{4. $C_{L,\max}$ est difficile pour tous.} Même avec 300 à 500 profils, l'erreur relative reste de 0,5 à 0,6 pour tous les modèles. La grandeur est irrégulière : les échecs d'XFOIL sur la rampe de manœuvre et la façon dont on retient le plus grand $C_L$ convergé créent des sauts d'un profil à son voisin. Aucun modèle lisse ne peut les représenter, et c'est sur cette grandeur que se joue la faisabilité du front.

\section{Sur le problème mis à l'échelle}
On refait la même étude sur les profils du problème CST (chapitre~\ref{ch:echelle}), en prenant pour données les profils simulés pendant la première campagne. Le modèle en $z$ a maintenant \textbf{66 coefficients} par grandeur ; la machine de factorisation en a 151 et le modèle bayésien 466.
\begin{figure}[htbp]
\centering
% figure générée par tools/make_figures.py
\begin{tikzpicture}
\begin{groupplot}[group style={group size=3 by 2,horizontal sep=1.15cm,vertical sep=2.3cm},width=0.33\linewidth,height=4.6cm,
  xmode=log,xlabel={profils d'entraînement $N$},tick label style={font=\scriptsize},label style={font=\scriptsize},
  title style={font=\small},legend style={font=\scriptsize,draw=none,at={(-0.75,-0.45)},anchor=north,legend columns=4}]
\nextgroupplot[title={traînée $C_D$},ylabel={erreur relative},ymin=0,ymax=1.2]
\addplot[cblue,very thick,mark=*,mark size=1.2pt] coordinates {(10,1.6496) (15,1.5445) (20,1.2272) (30,1.2371) (50,1.5387) (80,1.7084) (120,0.96605) (200,0.79414) (300,0.72964) (500,0.67401)};
\addplot[corange,very thick,mark=square*,mark size=1.2pt] coordinates {(10,1.0399) (15,1.038) (20,1.0528) (30,1.0062) (50,1.0604) (80,1.0624) (120,1.0861) (200,1.0678) (300,0.93575) (500,0.7612)};
\addplot[cgreen,very thick,mark=triangle*,mark size=1.2pt] coordinates {(10,1.0051) (15,0.99211) (20,0.97575) (30,0.93924) (50,0.89685) (80,0.88863) (120,0.82813) (200,0.78753) (300,0.76202) (500,0.73489)};
\addplot[cviolet,very thick,mark=diamond*,mark size=1.2pt] coordinates {(10,1.3268) (15,1.1865) (20,0.99513) (30,1.0322) (50,0.88074) (80,0.81249) (120,0.78326) (200,0.80308) (300,0.72709) (500,0.62105)};
\nextgroupplot[title={SER},ylabel={erreur relative},ymin=0,ymax=1.2]
\addplot[cblue,very thick,mark=*,mark size=1.2pt] coordinates {(10,1.2119) (15,1.0276) (20,0.91884) (30,0.73046) (50,0.56012) (80,0.17266) (120,0.12144) (200,0.089704) (300,0.085095) (500,0.078537)};
\addplot[corange,very thick,mark=square*,mark size=1.2pt] coordinates {(10,0.89888) (15,0.80702) (20,0.75167) (30,0.66848) (50,0.52418) (80,0.41089) (120,0.29845) (200,0.20303) (300,0.15222) (500,0.11385)};
\addplot[cgreen,very thick,mark=triangle*,mark size=1.2pt] coordinates {(10,0.88319) (15,0.79761) (20,0.71858) (30,0.58009) (50,0.38934) (80,0.31607) (120,0.27835) (200,0.25215) (300,0.23134) (500,0.1869)};
\addplot[cviolet,very thick,mark=diamond*,mark size=1.2pt] coordinates {(10,0.64754) (15,0.39075) (20,0.27033) (30,0.21768) (50,0.11191) (80,0.084752) (120,0.077736) (200,0.065246) (300,0.060054) (500,0.054701)};
\nextgroupplot[title={$C_{L,\max}$},ylabel={erreur relative},ymin=0,ymax=1.2]
\addplot[cblue,very thick,mark=*,mark size=1.2pt] coordinates {(10,2.147) (15,2.03) (20,1.9631) (30,1.9226) (50,2.0626) (80,1.4347) (120,0.97656) (200,0.75518) (300,0.6944) (500,0.65612)};
\addplot[corange,very thick,mark=square*,mark size=1.2pt] coordinates {(10,1.0214) (15,1.1017) (20,0.95844) (30,1.1037) (50,1.069) (80,1.0301) (120,1.1879) (200,1.0485) (300,0.8969) (500,0.71712)};
\addplot[cgreen,very thick,mark=triangle*,mark size=1.2pt] coordinates {(10,0.99577) (15,1.0169) (20,0.92292) (30,0.92755) (50,0.85232) (80,0.80111) (120,0.77597) (200,0.73498) (300,0.70434) (500,0.68611)};
\addplot[cviolet,very thick,mark=diamond*,mark size=1.2pt] coordinates {(10,1.2654) (15,1.2877) (20,1.0455) (30,0.98396) (50,0.95736) (80,0.9036) (120,0.894) (200,0.90686) (300,0.78768) (500,0.73796)};
\nextgroupplot[title={traînée $C_D$},ylabel={corrélation de rang},ymin=0,ymax=1]
\addplot[cblue,very thick,mark=*,mark size=1.2pt] coordinates {(10,0.19589) (15,0.26361) (20,0.3362) (30,0.37082) (50,0.33741) (80,0.37069) (120,0.64616) (200,0.73611) (300,0.76313) (500,0.80966)};
\addplot[corange,very thick,mark=square*,mark size=1.2pt] coordinates {(10,0.24038) (15,0.32291) (20,0.35623) (30,0.41448) (50,0.40841) (80,0.46494) (120,0.48334) (200,0.52989) (300,0.62815) (500,0.73671)};
\addplot[cgreen,very thick,mark=triangle*,mark size=1.2pt] coordinates {(10,0.24993) (15,0.31283) (20,0.36353) (30,0.45019) (50,0.51747) (80,0.53354) (120,0.6296) (200,0.67017) (300,0.70113) (500,0.73329)};
\addplot[cviolet,very thick,mark=diamond*,mark size=1.2pt] coordinates {(10,0.20575) (15,0.46828) (20,0.54958) (30,0.54888) (50,0.63792) (80,0.70763) (120,0.73864) (200,0.71668) (300,0.77647) (500,0.84335)};
\nextgroupplot[title={SER},ylabel={corrélation de rang},ymin=0,ymax=1]
\addplot[cblue,very thick,mark=*,mark size=1.2pt] coordinates {(10,0.61693) (15,0.63385) (20,0.73689) (30,0.78636) (50,0.88133) (80,0.98659) (120,0.99282) (200,0.99581) (300,0.99656) (500,0.99716)};
\addplot[corange,very thick,mark=square*,mark size=1.2pt] coordinates {(10,0.49401) (15,0.63244) (20,0.695) (30,0.79269) (50,0.88964) (80,0.93776) (120,0.96102) (200,0.97949) (300,0.98849) (500,0.99357)};
\addplot[cgreen,very thick,mark=triangle*,mark size=1.2pt] coordinates {(10,0.5869) (15,0.657) (20,0.75183) (30,0.84541) (50,0.93159) (80,0.95654) (120,0.95968) (200,0.96802) (300,0.97362) (500,0.98164)};
\addplot[cviolet,very thick,mark=diamond*,mark size=1.2pt] coordinates {(10,0.7724) (15,0.91659) (20,0.97142) (30,0.98021) (50,0.99422) (80,0.99655) (120,0.99699) (200,0.99783) (300,0.99831) (500,0.99853)};
\nextgroupplot[title={$C_{L,\max}$},ylabel={corrélation de rang},ymin=0,ymax=1]
\addplot[cblue,very thick,mark=*,mark size=1.2pt] coordinates {(10,-0.20489) (15,-0.20562) (20,-0.10526) (30,-0.050621) (50,0.18674) (80,0.43273) (120,0.59003) (200,0.70027) (300,0.75991) (500,0.7846)};
\addplot[corange,very thick,mark=square*,mark size=1.2pt] coordinates {(10,0.27866) (15,0.25074) (20,0.38232) (30,0.34959) (50,0.47324) (80,0.49835) (120,0.41627) (200,0.55481) (300,0.629) (500,0.73841)};
\addplot[cgreen,very thick,mark=triangle*,mark size=1.2pt] coordinates {(10,0.28108) (15,0.31317) (20,0.41874) (30,0.43034) (50,0.57438) (80,0.64976) (120,0.67772) (200,0.71466) (300,0.74959) (500,0.7661)};
\addplot[cviolet,very thick,mark=diamond*,mark size=1.2pt] coordinates {(10,0.3568) (15,0.40905) (20,0.45653) (30,0.51508) (50,0.55018) (80,0.57532) (120,0.5778) (200,0.60261) (300,0.69381) (500,0.7457)};
\legend{quadratique en $z$, machine de factorisation, bayésien sur les bits, réseau de neurones}
\end{groupplot}
\end{tikzpicture}

\caption{Précision des quatre modèles sur le problème CST (30 bits), même présentation que la figure~\ref{fig:precision-v2}.}\label{fig:precision-cst}
\end{figure}
\begin{center}
\renewcommand{\arraystretch}{1.2}
\small
\begin{tabular}{l ccc ccc}
\toprule
& \multicolumn{3}{c}{erreur relative} & \multicolumn{3}{c}{corrélation de rang} \\
SER, $N=$ & 30 & 80 & 300 & 30 & 80 & 300 \\
\midrule
quadratique en $z$ & 0,73 & 0,17 & 0,09 & 0,79 & 0,99 & 1,00 \\
machine de factorisation & 0,67 & 0,41 & 0,15 & 0,79 & 0,94 & 0,99 \\
bayésien sur les bits & 0,58 & 0,32 & 0,23 & 0,85 & 0,96 & 0,97 \\
réseau de neurones & 0,22 & 0,08 & 0,06 & 0,98 & 1,00 & 1,00 \\
\midrule
$C_D$, $N=$ & 30 & 80 & 300 & 30 & 80 & 300 \\
\midrule
quadratique en $z$ & 1,24 & \textbf{1,71} & 0,73 & 0,37 & 0,37 & 0,76 \\
machine de factorisation & 1,01 & 1,06 & 0,94 & 0,41 & 0,46 & 0,63 \\
bayésien sur les bits & 0,94 & 0,89 & 0,76 & 0,45 & 0,53 & 0,70 \\
réseau de neurones & 1,03 & 0,81 & 0,73 & 0,55 & 0,71 & 0,78 \\
\midrule
$C_{L,\max}$, $N=$ & 30 & 80 & 300 & 30 & 80 & 300 \\
\midrule
quadratique en $z$ & 1,92 & 1,43 & 0,69 & $-0{,}05$ & 0,43 & 0,76 \\
machine de factorisation & 1,10 & 1,03 & 0,90 & 0,35 & 0,50 & 0,63 \\
bayésien sur les bits & 0,93 & 0,80 & 0,70 & 0,43 & 0,65 & 0,75 \\
réseau de neurones & 0,98 & 0,90 & 0,79 & 0,52 & 0,58 & 0,69 \\
\bottomrule
\end{tabular}
\end{center}

\paragraph{5. Le pic « juste déterminé » se déplace avec le nombre de coefficients.} Sur $C_D$, l'erreur du modèle en $z$ culmine à 1,71 pour $N=80$, juste au-dessus de ses 66 coefficients, comme elle culminait à $N=15$ en V2 (15 coefficients). Tant que $N$ n'a pas largement dépassé 66, le modèle en $z$ classe mal la traînée et la contrainte : sa corrélation de rang reste sous 0,45 jusqu'à 80 profils. Les modèles en bits, régularisés (pénalité $L_2$ pour la machine de factorisation, a priori bayésien pour BOCS), n'ont pas ce pic. \textbf{Le rapport s'est inversé} : en V2, la boucle disposait de 5 fois plus de profils que de coefficients du modèle en $z$ ; à l'échelle, elle en a à peine autant.

\paragraph{6. La traînée devient la grandeur difficile.} Même à 300 profils, aucun modèle ne descend sous 0,73 d'erreur relative sur $C_D$, alors que la SER est très bien prédite (0,06 à 0,23). Deux explications plausibles, non vérifiées séparément : la traînée calculée par XFOIL varie fortement près du décollement, et les profils de la base, choisis par les méthodes d'optimisation, sont concentrés dans les régions extrêmes de la grille, où c'est le cas.

\begin{piege}[ : précision moyenne et utilité dans la boucle]
Le chapitre suivant montrera que le modèle le plus précis en moyenne, le réseau de neurones, est le \emph{moins} bon dans la boucle. La précision moyenne sur toute la grille n'est pas ce qui compte : ce qui compte est de bien classer les profils \emph{près du front}, et de ne pas exclure à tort les profils du front par la contrainte. Une mesure de précision globale ne le dit pas.
\end{piege}

\begin{resultat}
\begin{itemize}
\item Le réseau de neurones est le plus précis, même avec 10 à 30 profils.
\item Les modèles en bits ne rattrapent le modèle en $z$ qu'au-delà de 80 à 120 profils.
\item Le modèle en $z$ est fragile autour de $N=15$, nombre de ses coefficients.
\item $C_{L,\max}$ reste mal prédit par tous (erreur relative de 0,5 à 0,6) : la contrainte est la grandeur la plus incertaine.
\item À l'échelle, le modèle en $z$ a 66 coefficients, et son pic de mauvaise précision arrive vers 80 profils, au milieu du budget : les modèles en bits régularisés prennent l'avantage.
\end{itemize}
\end{resultat}

\exercices
\begin{exercice}\label{ex:prec-lecture}
Un modèle a une erreur relative de 0,35 sur $C_D$, dont l'écart-type sur la grille vaut environ 10 counts. Quelle est son erreur en counts ? Les profils du front de S2 vont de 38 à 59 counts (hors profil aberrant) : cette erreur suffit-elle à les classer correctement ?
\end{exercice}
\begin{exercice}\label{ex:prec-rang}
Pourquoi la corrélation de rang peut-elle être bonne (0,93) alors que l'erreur relative est médiocre (0,35) ? Donner un exemple de modèle qui aurait une corrélation de rang de 1 et une erreur relative de 1.
\end{exercice}

\begin{pistes}
\begin{itemize}
\item Mesurer la précision seulement sur les profils proches du front, ou pondérée par la proximité au front, plutôt que sur toute la grille.
\item Régulariser le modèle en $z$ (ridge choisi par validation croisée au lieu de $\lambda=10^{-6}$), ce qui devrait effacer son pic à $N\approx66$.
\end{itemize}
\end{pistes}

\chapter{Les contraintes : modélisées ou laissées à la boîte noire}\label{ch:contraintes}
\iterbanner{V3}

\section{Deux façons de faire}
Le profil doit être assez épais ($t\ge8~\%$), tenir la manœuvre ($C_{L,\max}\ge1{,}2$) et faire converger XFOIL en croisière. La boucle de la V2 traite ces contraintes par des pénalités dans le QUBO (cours de calcul quantique, chapitre~8). Ce cours compare deux approches.

\paragraph{Contrainte modélisée (\code{modele}), comme en V2.}
\begin{itemize}
\item l'épaisseur est connue sans simulation : l'indice $x_t=0$ est interdit par une pénalité exacte, rendue quadratique par un qubit auxiliaire (Rosenberg) ;
\item $C_{L,\max}$ est \emph{modélisé} par un modèle affine en $z$, ajusté sur les profils où XFOIL a convergé, et pénalisé avec 3 bits d'écart (poids 75) ;
\item la traînée et la SER sont modélisées sur les seuls profils faisables.
\end{itemize}
Le QUBO compte 14 qubits sur le problème V2 (10 de conception, 1 auxiliaire, 3 d'écart), et 33 sur le problème mis à l'échelle (30 de conception, 3 d'écart, pas d'auxiliaire car l'épaisseur n'y est plus une variable).

\paragraph{Contrainte en boîte noire (\code{bn}), comme FMQA.} On ne modélise aucune contrainte. Pour chaque poids $w$, on construit directement les données de l'objectif pondéré :
\[ y_n=\begin{cases}w\,C_{D,n}/s_{C_D}+(1-w)\,{\rm SER}_n/s_{\rm SER} & \text{si le profil }n\text{ est faisable,}\\[2pt] \max_{\text{faisables}}y+0{,}2\times\big(\max_{\text{faisables}}y-\min_{\text{faisables}}y\big) & \text{sinon,}\end{cases} \]
et l'on ajuste un seul modèle sur \emph{tous} les profils simulés. Le modèle apprend que les régions infaisables sont mauvaises, comme il apprend que les régions traînantes le sont. Seule l'épaisseur, connue sans calcul, garde sa pénalité exacte. Le QUBO n'a plus de bits d'écart : 11 qubits sur le problème V2, 30 à l'échelle.

\paragraph{Exemple.} Si les profils faisables ont des objectifs pondérés entre 0,2 et 1,0, un profil infaisable reçoit $1{,}0+0{,}2\times0{,}8=1{,}16$ : un peu plus mauvais que le plus mauvais faisable. Le choix de 0,2 est arbitraire ; trop grand, il crée une marche que le modèle quadratique lisse mal ; trop petit, l'infaisable n'est plus assez repoussé.

\section{Avantages et inconvénients}
\begin{center}
\renewcommand{\arraystretch}{1.25}
\small
\begin{tabularx}{\linewidth}{>{\raggedright\arraybackslash}X>{\raggedright\arraybackslash}X}
\toprule
Contrainte modélisée & Contrainte en boîte noire \\
\midrule
utilise l'information « de combien la contrainte est violée » ($C_{L,\max}=1{,}1$ n'est pas $0{,}6$) & ne sait que « faisable ou non » : l'écart à la contrainte est perdu \\
modèle de $C_{L,\max}$ forcé affine (sinon le carré est de degré 4) & aucun modèle de contrainte, donc aucune forme imposée \\
3 bits d'écart, pénalité en dents de scie qui domine le spectre & pas de bits d'écart ; l'objectif occupe une plus grande part du spectre \\
2 modèles par itération (traînée, SER) & 5 modèles par itération (un par poids) \\
la frontière de faisabilité est une surface lisse prédite & la frontière est une marche, que les modèles lisses représentent mal \\
\bottomrule
\end{tabularx}
\end{center}

\section{Le QUBO en boîte noire est-il plus facile pour QAOA ?}
Le cours de calcul quantique (chapitre~15) a montré que, dans le QUBO de la V2, l'objectif n'occupe que 1,7~\% de l'étendue des énergies : les pénalités écrasent tout, et QAOA ne fait que 5 fois mieux que le hasard. Sans bits d'écart, la situation devrait s'améliorer. On le mesure (\code{ia\_qubo/qaoa\_contraintes.py}) sur 9 QUBO de début de boucle (3 graines, 3 poids), avec le modèle en $z$ dans les deux cas pour n'isoler que l'effet de la contrainte :
\begin{center}
\renewcommand{\arraystretch}{1.2}
\begin{tabularx}{\linewidth}{>{\raggedright\arraybackslash}X c >{\centering\arraybackslash}p{3.3cm} >{\centering\arraybackslash}p{3.3cm}}
\toprule
& qubits & part de l'objectif dans le spectre & QAOA $p=2$ : gain sur le hasard \\
\midrule
contrainte modélisée & 14 & 1,7~\% & 4,5 (3,5 à 4,8) \\
contrainte en boîte noire & 11 & 9,8~\% & 4,1 (0,1 à 8,3) \\
\bottomrule
\end{tabularx}
\end{center}
(Médianes, entre parenthèses le minimum et le maximum.) La part de l'objectif est multipliée par six, mais QAOA n'en profite pas en médiane. La raison est visible dans la matrice $Q$ : la pénalité d'épaisseur reste, avec ses termes de Rosenberg (couplages jusqu'à 12), qui dominent encore un objectif normalisé d'étendue de l'ordre de 1. Supprimer \emph{une} pénalité ne suffit pas : il faudrait qu'aucune ne domine. C'est l'argument pour les espaces latents qui ne décrivent que des conceptions valides (chapitre~\ref{ch:ia-autres}), ou pour des mélangeurs qui respectent les contraintes (cours de calcul quantique, chapitre~15).

\begin{piege}[ : la valeur dégradée des infaisables]
Dans la boîte noire, la valeur donnée aux infaisables est \emph{inventée}. Un modèle quadratique ajusté sur des données qui sautent brusquement à la frontière de faisabilité fait des erreurs des deux côtés : il peut rendre attractif un profil infaisable voisin d'un bon faisable, ou repousser un bon profil faisable voisin d'un infaisable. Or le front est justement au bord de la contrainte de manœuvre (lot L1).
\end{piege}

\begin{resultat}
\begin{itemize}
\item Contrainte modélisée : plus d'information (l'écart à la contrainte), mais un modèle affine imposé et des bits d'écart qui dominent le spectre.
\item Contrainte en boîte noire : aucun modèle de contrainte, un QUBO plus petit où l'objectif pèse six fois plus, mais une frontière apprise en marche d'escalier.
\item Pour QAOA, retirer les bits d'écart ne suffit pas : la pénalité d'épaisseur domine encore.
\end{itemize}
\end{resultat}

\exercices
\begin{exercice}\label{ex:bn-valeur}
Dix profils simulés : six faisables d'objectifs pondérés $0{,}3;\ 0{,}5;\ 0{,}6;\ 0{,}8;\ 0{,}9;\ 1{,}1$, et quatre infaisables. Quelle valeur reçoivent les infaisables ? Que se passe-t-il si le prochain profil simulé est faisable avec un objectif de $1{,}5$ ?
\end{exercice}
\begin{exercice}\label{ex:bn-qubits}
Combien de qubits et de couplages entre bits de conception et bits d'écart économise la boîte noire sur le problème V2 ? Sur le problème mis à l'échelle (30 bits de conception) ?
\end{exercice}

\begin{pistes}
\begin{itemize}
\item Modéliser la faisabilité par un \emph{classifieur} (probabilité d'être faisable) plutôt que par une valeur dégradée, et l'utiliser comme dans l'optimisation bayésienne contrainte du lot L1.
\item Retirer aussi la pénalité d'épaisseur en décalant la grille de $t$ (8 niveaux de 8 à 22~\% au lieu de 6 à 20~\%) : la contrainte disparaît de la formulation, sans auxiliaire ni pénalité. La formulation du problème est un levier plus puissant que le solveur.
\end{itemize}
\end{pistes}

\chapter{Résultats sur le problème V2}\label{ch:resultats-v2}
\iterbanner{V3}

\section{Le protocole}
Huit méthodes (4 modèles $\times$ 2 traitements de la contrainte), deux scénarios radar (S1, S2), 20 graines chacune, avec le protocole du lot L1 : mêmes 10 profils initiaux pour une graine donnée, 5 poids par itération, budget de 250 simulations, SER par méthode des moments, QUBO résolus exactement par énumération. La campagne a pris 46 minutes sur 11 processus (\code{python3 -m ia\_qubo.experiences\_v2\_ia}). La méthode \code{quadz\_modele} redonne exactement les chiffres de la V2 (75 et 142,5) : c'est attendu, puisque la boucle générique la reproduit à l'identique.

\section{Le front exact : qui le retrouve, et en combien de simulations}
\begin{center}
\renewcommand{\arraystretch}{1.2}
\small
\begin{tabular}{lcccccc}
\toprule
& \multicolumn{3}{c}{S1 (front de 4 profils)} & \multicolumn{3}{c}{S2 (front de 10 profils)} \\
Méthode & médiane & quartiles & réussite & médiane & quartiles & réussite \\
\midrule
\textbf{quadratique en $z$ (actuel)} & \textbf{75} & 60 -- 96 & 100~\% & \textbf{142,5} & 116 -- 191 & 100~\% \\
machine de factorisation & 95 & 75 -- 132 & 100~\% & 152,5 & 130 -- 211 & 95~\% \\
bayésien sur les bits & 100 & 90 -- 115 & 100~\% & 187,5 & 159 -- 245 & 80~\% \\
réseau de neurones & --- & & 5~\% & --- & & 5~\% \\
\midrule
les quatre, en boîte noire & --- & & 0~\% & --- & & 0~\% \\
\bottomrule
\end{tabular}
\end{center}
(Nombre de simulations pour retrouver tout le front exact ; « --- » : la médiane n'est pas atteinte dans le budget de 250, faute de réussite dans au moins la moitié des graines.)

Comparaisons appariées avec l'approche actuelle, graine par graine (test de Wilcoxon) :
\begin{itemize}
\item \textbf{machine de factorisation} : 22,5 (S1) et 17,5 (S2) simulations de plus en médiane, mais elle fait mieux sur 7 et 6 graines sur 20. L'écart est à la limite de la significativité ($p=0{,}052$ et $0{,}057$) ;
\item \textbf{modèle bayésien sur les bits} : significativement moins bon ($p=0{,}007$ et $0{,}015$) ;
\item \textbf{réseau de neurones} : très nettement moins bon ($p=10^{-4}$), alors qu'il est le plus précis (chapitre~\ref{ch:precision}).
\end{itemize}

\begin{figure}[htbp]
\centering
% figure générée par tools/make_figures.py
\begin{tikzpicture}
\begin{axis}[width=0.95\linewidth,height=6.2cm,xlabel={profils simulés},ylabel={part du front exact retrouvée (médiane)},
  xmin=10,xmax=250,ymin=0,ymax=1.02,tick label style={font=\scriptsize},label style={font=\small},
  legend style={font=\scriptsize,draw=none,at={(0.5,-0.22)},anchor=north,legend columns=2},legend cell align=left]
\addplot[cblue,very thick,mark=*,mark size=1pt,mark repeat=3] coordinates {(10,0) (20,0) (30,0.1) (40,0.1) (50,0.3) (60,0.4) (70,0.55) (80,0.7) (90,0.75) (100,0.85) (110,0.85) (120,0.9) (130,0.9) (140,0.95) (150,1) (160,1) (170,1) (180,1) (190,1) (200,1) (210,1) (220,1) (230,1) (240,1) (250,1)};
\addplot[corange,very thick,mark=square*,mark size=1pt,mark repeat=3] coordinates {(10,0) (20,0) (30,0.1) (40,0.2) (50,0.3) (60,0.4) (70,0.5) (80,0.6) (90,0.7) (100,0.8) (110,0.8) (120,0.85) (130,0.9) (140,0.9) (150,0.95) (160,1) (170,1) (180,1) (190,1) (200,1) (210,1) (220,1) (230,1) (240,1) (250,1)};
\addplot[cgreen,very thick,mark=triangle*,mark size=1pt,mark repeat=3] coordinates {(10,0) (20,0) (30,0.1) (40,0.1) (50,0.2) (60,0.2) (70,0.3) (80,0.4) (90,0.5) (100,0.55) (110,0.6) (120,0.7) (130,0.75) (140,0.8) (150,0.85) (160,0.9) (170,0.9) (180,0.9) (190,1) (200,1) (210,1) (220,1) (230,1) (240,1) (250,1)};
\addplot[cviolet,very thick,mark=diamond*,mark size=1pt,mark repeat=3] coordinates {(10,0) (20,0) (30,0.1) (40,0.3) (50,0.4) (60,0.5) (70,0.7) (80,0.85) (90,0.9) (100,0.9) (110,0.9) (120,0.9) (130,0.9) (140,0.9) (150,0.9) (160,0.9) (170,0.9) (180,0.9) (190,0.9) (200,0.9) (210,0.9) (220,0.9) (230,0.9) (240,0.9) (250,0.9)};
\addplot[cblue,very thick,mark=*,dashed,mark size=1pt,mark repeat=3] coordinates {(10,0) (20,0) (30,0) (40,0) (50,0) (60,0) (70,0) (80,0) (90,0) (100,0) (110,0) (120,0.05) (130,0.05) (140,0.05) (150,0.05) (160,0.05) (170,0.15) (180,0.2) (190,0.25) (200,0.35) (210,0.4) (220,0.4) (230,0.4) (240,0.4) (250,0.4)};
\addplot[corange,very thick,mark=square*,dashed,mark size=1pt,mark repeat=3] coordinates {(10,0) (20,0) (30,0) (40,0.05) (50,0.1) (60,0.1) (70,0.1) (80,0.1) (90,0.1) (100,0.1) (110,0.15) (120,0.2) (130,0.2) (140,0.2) (150,0.2) (160,0.25) (170,0.25) (180,0.25) (190,0.25) (200,0.3) (210,0.3) (220,0.3) (230,0.3) (240,0.3) (250,0.35)};
\addplot[cgreen,very thick,mark=triangle*,dashed,mark size=1pt,mark repeat=3] coordinates {(10,0) (20,0) (30,0) (40,0) (50,0) (60,0.05) (70,0.05) (80,0.05) (90,0.1) (100,0.15) (110,0.2) (120,0.2) (130,0.2) (140,0.3) (150,0.3) (160,0.35) (170,0.4) (180,0.4) (190,0.5) (200,0.55) (210,0.6) (220,0.6) (230,0.6) (240,0.6) (250,0.6)};
\addplot[cviolet,very thick,mark=diamond*,dashed,mark size=1pt,mark repeat=3] coordinates {(10,0) (20,0) (30,0) (40,0) (50,0) (60,0) (70,0.05) (80,0.1) (90,0.15) (100,0.2) (110,0.25) (120,0.3) (130,0.3) (140,0.35) (150,0.35) (160,0.35) (170,0.35) (180,0.4) (190,0.4) (200,0.4) (210,0.4) (220,0.4) (230,0.4) (240,0.4) (250,0.4)};
\addplot[cteal,very thick,mark=pentagon*,mark size=1pt,mark repeat=3] coordinates {(10,0) (20,0) (30,0.1) (40,0.25) (50,0.4) (60,0.7) (70,0.85) (80,0.9) (90,1) (100,1) (110,1) (120,1) (130,1) (140,1) (150,1) (160,1) (170,1) (180,1) (190,1) (200,1) (210,1) (220,1) (230,1) (240,1) (250,1)};
\legend{{quadratique en $z$ (actuel)}, {machine de factorisation}, {bayésien sur les bits}, {réseau de neurones}, {quadratique en $z$, boîte noire}, {factorisation, boîte noire}, {bayésien, boîte noire}, {réseau, boîte noire}, {hybride (réseau + contrainte affine)}}
\end{axis}
\end{tikzpicture}

\caption{Scénario S2 : part du front exact retrouvée en fonction du nombre de profils simulés (médiane sur 20 graines). En trait plein, contrainte modélisée ; en tirets, contrainte en boîte noire.}\label{fig:v2-s2}
\end{figure}

\section{Ce qui s'est passé}
\paragraph{1. Parmi les modèles qui sont des QUBO, le plus contraint gagne.} Pour les trois modèles quadratiques, le classement suit le nombre de paramètres : 15 (en $z$), 51 (factorisation), 56 (bayésien). Les modèles en bits ne rejoignent la précision du modèle en $z$ qu'au-delà de 80 à 120 profils (chapitre~\ref{ch:precision}) ; or la boucle trouve le front de S1 en 75 simulations, dont une partie infaisable. Le modèle en $z$ apporte une information que les données ne donnent pas : la continuité des niveaux (chapitre~\ref{ch:fm}, piège). Le réseau de neurones fait exception : il est le plus précis, même avec peu de données, et pourtant il échoue. La raison est ailleurs.

\paragraph{2. Le réseau de neurones échoue à cause de la contrainte, pas de l'objectif.} En fin de boucle, il trouve 90~\% du front de S2 et la moitié de celui de S1, mais il manque toujours les mêmes profils : des profils à nez fin ($I=2$) et d'épaisseur 10 à 14~\%. On a d'abord pensé à des points non supportés (aucune somme pondérée ne les désigne), et l'on a vérifié (\code{ia\_qubo/diagnostic\_v2\_ia.py}) : l'hypothèse est \emph{fausse}, le réseau manque aussi des points supportés. Ce qu'il prédit pour ces profils, après 250 simulations :
\begin{center}
\renewcommand{\arraystretch}{1.2}
\small
\begin{tabular}{lccc}
\toprule
Profil du front de S2 & $C_{L,\max}$ vrai & prédit par le réseau & prédit par le modèle affine (V2) \\
\midrule
(5~\%, 0,45, 10~\%, $I=2$) & 1,34 & \textbf{0,89} & 1,30 \\
(4~\%, 0,55, 12~\%, $I=2$) & 1,20 & \textbf{0,96} & 1,24 \\
\bottomrule
\end{tabular}
\end{center}
Le réseau les croit nettement infaisables, et ne les simule donc jamais. Pourquoi ? Autour d'eux, les profils voisins à nez fin ont des $C_{L,\max}$ de 0,9 à 1,1, ou font échouer XFOIL sur la rampe de manœuvre : la grandeur est très irrégulière dans ce coin de la grille. Le réseau, précis et souple, suit fidèlement ces voisins. Le modèle affine, trop grossier pour les suivre, fait une moyenne indulgente qui laisse passer les bons profils. C'est exactement la conclusion du lot L1 : \textbf{le front est au bord de la contrainte, et un modèle de contrainte plus fidèle est plus sévère, donc il l'exclut}.

\paragraph{3. La boîte noire échoue pour la même raison, en pire.} Les variantes en boîte noire ne retrouvent jamais tout le front : au mieux 60~\% de celui de S2 (modèle bayésien, médiane), souvent beaucoup moins. Le front est collé à la frontière de faisabilité ; or c'est là que les données sautent de la vraie valeur à la valeur dégradée des infaisables, et un modèle lisse étale ce saut sur les deux côtés. Les profils du front, voisins d'infaisables, sont jugés mauvais. L'hypervolume atteint pourtant 0,91 à 0,98 de celui du front exact en fin de budget : la boîte noire trouve de bons compromis, mais pas les meilleurs, qui sont au bord.

\section{Le temps de calcul}
Par graine, en moyenne : 1 à 3 secondes pour le modèle en $z$ et le modèle bayésien, 18 à 31 secondes pour la machine de factorisation, 3 minutes pour le réseau de neurones (contrainte modélisée) et 6 à 7 minutes en boîte noire (un réseau par poids). Comparé à 250 simulations XFOIL et MoM, tout cela reste négligeable : ce n'est pas le coût d'apprentissage qui départage les méthodes.

\begin{resultat}
\begin{itemize}
\item Sur le problème V2, l'approche actuelle reste la meilleure. La machine de factorisation est proche (écart non significatif à 5~\%), le modèle bayésien sur les bits et surtout le réseau de neurones sont moins bons.
\item Plus un modèle est souple, plus il lui faut de données ; à ce budget, la continuité imposée par le modèle en $z$ vaut plus que la souplesse.
\item Le modèle le plus précis échoue parce qu'il modélise trop fidèlement une contrainte irrégulière, et exclut le front, qui est à son bord.
\item La contrainte en boîte noire ne retrouve jamais tout le front, pour la même raison.
\end{itemize}
\end{resultat}

\exercices
\begin{exercice}\label{ex:v2-lecture}
D'après le tableau, combien de simulations la machine de factorisation coûte-t-elle de plus que l'approche actuelle dans S1, en médiane ? Pourquoi la comparaison graine par graine (test apparié) est-elle plus pertinente que la comparaison des deux médianes ?
\end{exercice}
\begin{exercice}\label{ex:v2-contrainte}
Un profil a un $C_{L,\max}$ vrai de 1,34. Le réseau prédit 0,89. Quelle pénalité lui ajoute la boucle ($75\max(0,1{,}2-\widehat{C_{L,\max}})^2$) ? Comparer à l'étendue de l'objectif normalisé (environ 1). Que faudrait-il pour que ce profil soit quand même proposé ?
\end{exercice}

\begin{pistes}
\begin{itemize}
\item Garder le réseau de neurones pour la traînée et la SER, mais un modèle affine (ou un classifieur indulgent) pour la contrainte : cela séparerait l'effet du modèle d'objectif de celui du modèle de contrainte.
\item Utiliser l'écart entre les membres de l'ensemble (incertitude) pour relâcher la contrainte là où le réseau doute, comme la probabilité de faisabilité du lot L1.
\end{itemize}
\end{pistes}

\chapter{Un problème mis à l'échelle}\label{ch:echelle}
\iterbanner{V3}

\section{Pourquoi un second problème}
Le problème V2 a deux propriétés qui le rendent confortable, mais irréaliste : on peut \textbf{tout énumérer} (1\,024 profils, $2^{14}$ chaînes du QUBO), et l'on connaît le \textbf{front exact}. Sur un vrai problème de conception, aucune des deux n'est vraie. Pour savoir ce que deviennent les méthodes quand on s'en approche, on construit un second problème :
\begin{itemize}
\item une forme décrite par \textbf{10 paramètres} au lieu de 4 ;
\item \textbf{30 bits} : $8^{10}\approx1{,}07\cdot10^9$ profils, qu'on ne peut évidemment pas tous simuler, et un QUBO de 30 à 33 qubits, dont l'énumération prendrait de l'ordre d'une minute par QUBO, pour environ 300 QUBO par exécution : il faut un solveur heuristique, le recuit simulé ;
\item des \textbf{solveurs boîtes noires} appelés à la demande (XFOIL, SER), avec un coût réel par évaluation ;
\item un front \textbf{inconnu}.
\end{itemize}

\section{La forme : la paramétrisation CST}
La méthode CST (\emph{Class-Shape Transformation}, Kulfan, 2008) est très utilisée en conception aéronautique. Chaque face du profil s'écrit comme le produit d'une \emph{fonction de classe}, qui donne l'allure générale, et d'une \emph{fonction de forme}, qui la module :
\[ y(x)=\underbrace{\sqrt x\,(1-x)}_{C(x)}\ \underbrace{\sum_{i=0}^{4}A_i\,\binom4i\,x^i(1-x)^{4-i}}_{S(x)}. \]
\begin{itemize}
\item $C(x)=\sqrt x(1-x)$ donne un nez arrondi ($y\propto\sqrt x$ près de $x=0$) et un bord de fuite pointu ($y\to0$ en $x=1$).
\item $S(x)$ est une combinaison de 5 polynômes de Bernstein de degré 4. Chaque coefficient $A_i$ agit surtout sur une portion de la corde : $A_0$ au nez, $A_4$ près du bord de fuite.
\item Le rayon de bord d'attaque de chaque face vaut $A_0^2/2$ : le premier coefficient pilote directement le nez, donc la SER.
\end{itemize}
Avec 5 coefficients pour l'extrados ($A_{e,0..4}$, positifs) et 5 pour l'intrados ($A_{i,0..4}$, plutôt négatifs), on a 10 variables.

\paragraph{Exemple : le NACA 2412 en CST.} En ajustant les coefficients par moindres carrés sur le contour du NACA 2412, on trouve
\[ \mathbf A_e=(0{,}192;\ 0{,}213;\ 0{,}208;\ 0{,}194;\ 0{,}218),\qquad\mathbf A_i=(-0{,}152;\ -0{,}094;\ -0{,}100;\ -0{,}076;\ -0{,}078), \]
avec un écart maximal au contour de $1{,}4\cdot10^{-3}$ corde. Les rayons de nez valent $0{,}192^2/2=0{,}018$ (extrados) et $0{,}152^2/2=0{,}012$ (intrados), à comparer à $0{,}016$ pour la formule du NACA : la CST représente le profil avec 10 nombres, à la précision d'un tracé.

\paragraph{Les bornes.} Pour que la grille couvre des profils réalistes, on ajuste la CST sur 60 profils NACA (cambrure de 0 à 6~\%, position 0,3 à 0,5, épaisseur de 8 à 18~\%), on prend pour chaque coefficient le minimum et le maximum, élargis de 10~\%, et l'on y place 8 niveaux régulièrement espacés. Le NACA 2412 tombe près des indices $(2,3,2,3,2;\ 4,3,3,3,3)$. Sur 30 profils tirés au hasard dans la grille, 29 ont une géométrie valide et 23 sont faisables.

\begin{codebox}[ : \code{ia\_qubo/probleme\_cst.py}]
\begin{lstlisting}
def cst_coords(Au, Al, n_half=120):
    x = 0.5 * (1 - np.cos(np.linspace(0, np.pi, n_half + 1)))     # points resserres au nez et au BF
    yu = classe(x) * (bernstein_matrix(x) @ Au)                     # extrados
    yl = classe(x) * (bernstein_matrix(x) @ Al)                     # intrados
    return np.concatenate([x[::-1], x[1:]]), np.concatenate([yu[::-1], yl[1:]])
\end{lstlisting}
\end{codebox}

\section{Les solveurs, en boîtes noires}
Chaque évaluation d'un profil (\code{evaluer}) fait, dans l'ordre :
\begin{enumerate}
\item \textbf{la géométrie} : contour CST, vérification que l'intrados reste sous l'extrados, épaisseur maximale ;
\item \textbf{XFOIL en croisière} : $C_L=0{,}4$, $Re=3\cdot10^6$ $\to$ incidence et traînée $C_D$ ;
\item \textbf{XFOIL en manœuvre} : rampe de 4 à 14° $\to$ $C_{L,\max}$ ;
\item \textbf{la SER} : optique physique sur le secteur S2 ($-60^\circ$ à $0^\circ$, 81 directions, 10~GHz), à l'incidence de croisière.
\end{enumerate}
Contraintes : épaisseur maximale $\ge8~\%$, $C_{L,\max}\ge1{,}2$, convergence d'XFOIL, géométrie valide. Comme pour le profil aberrant de la V2, une traînée de plus de 200 counts à $C_L=0{,}4$ compte comme un échec d'XFOIL : 85 profils de la première campagne, jusqu'à 967 counts, sont des artefacts du solveur près du décollement. Une évaluation prend de 0,5 à 2 secondes.

\begin{piege}[ : pourquoi l'optique physique et pas la méthode des moments]
À l'échelle, chaque méthode demande 300 évaluations, et l'on compare 7 méthodes sur 8 graines : la première campagne a simulé 16\,242 profils distincts. La méthode des moments, à plusieurs secondes par profil et par secteur, rendrait l'étude trop longue. L'optique physique sous-estime la SER d'environ 1~dB mais garde le classement (corrélation de rang 0,999 sur la grille V2, cours SER) : c'est le classement qui compte pour comparer des méthodes d'optimisation. Sur un cas réel, ce serait exactement la place d'une approche multi-fidélité (chapitre~\ref{ch:ia-autres}).
\end{piege}

\paragraph{Le cache.} Chaque profil évalué est écrit dans \code{resultats/cache\_cst.jsonl}. Un profil demandé par deux méthodes n'est calculé qu'une fois ; mais le \emph{compteur} de chaque exécution compte toujours les profils distincts qu'elle a demandés, comme \code{CountedProblem}. Le cache économise du temps de calcul, pas des simulations au sens de la mesure.

\section{Mesurer sans connaître le front}
Le front exact est inconnu. On construit un \textbf{front de référence} : le front non dominé de tous les profils faisables évalués, par toutes les méthodes et toutes les graines. Pour chaque exécution, on calcule tous les 20 profils simulés l'hypervolume de ses profils faisables, divisé par celui du front de référence (objectifs normalisés comme en V2, point de référence $(1{,}1;\ 1{,}1)$). Une méthode qui atteint 1 a trouvé un front aussi bon que la réunion de tous les autres.

\section{Le solveur des QUBO : un recuit réglé}
Avec 30 à 33 qubits, énumérer les $2^{33}\approx8\cdot10^9$ chaînes de chaque QUBO n'est plus raisonnable dans la boucle : les QUBO sont résolus par recuit simulé. Le cours de calcul quantique (chapitre~9) a montré que le réglage de la boucle V2, avec une température proportionnelle au plus grand coefficient de $Q$ (donc aux pénalités), trouve rarement l'optimum. On a donc mesuré trois réglages sur 18 QUBO de la boucle V2 (6 graines, 3 poids), dont l'optimum est connu par énumération, avec 64 lectures :
\begin{center}
\renewcommand{\arraystretch}{1.2}
\begin{tabularx}{\linewidth}{>{\raggedright\arraybackslash}X >{\centering\arraybackslash}p{2.4cm} >{\centering\arraybackslash}p{3cm}}
\toprule
Température (décroissance géométrique) & lectures optimales & meilleure lecture $=$ optimum \\
\midrule
de $2\max|Q|$ à $0{,}01\max|Q|$ (réglage V2, pénalités) & 0,7~\% & 39~\% \\
de $2\max|Q|$ à $0{,}01\max|Q_{\rm obj}|$, 300 balayages & 1,4~\% & 56~\% \\
de $2\max|Q_{\rm obj}|$ à $0{,}01\max|Q_{\rm obj}|$ (objectif) & 2,7~\% & \textbf{67~\%} \\
\bottomrule
\end{tabularx}
\end{center}
Le réglage sur l'objectif est le meilleur : c'est lui qu'utilise \code{ia\_qubo/boucle.py}, avec 128 lectures et 200 balayages, pour toutes les méthodes à QUBO. La comparaison des modèles se fait donc à solveur égal. Le recuit est vectorisé (les 128 lectures avancent ensemble, en NumPy) : un QUBO de 33 qubits se résout en une fraction de seconde.

\begin{center}
\renewcommand{\arraystretch}{1.2}
\small
\begin{tabular}{lccc}
\toprule
Méthode & paramètres du modèle & qubits & choix du candidat \\
\midrule
hasard & --- & --- & tirage uniforme \\
NSGA-II (population 16) & --- & --- & algorithme génétique \\
\code{quadz\_modele} (approche actuelle) & 66 par grandeur & 33 & recuit \\
\code{fm\_modele} & 151 & 33 & recuit \\
\code{fm\_bn} & 151 (un par poids) & 30 & recuit \\
\code{bocs\_modele} & 466 & 33 & recuit (tirage de Thompson) \\
\code{hybride\_modele} (témoin) & $2\times1\,441+11$ & --- & tirage et recherche locale \\
\bottomrule
\end{tabular}
\end{center}

\begin{resultat}
\begin{itemize}
\item Le problème mis à l'échelle a 10 paramètres CST, 30 bits, $10^9$ profils, un front inconnu et des solveurs appelés à la demande.
\item La CST décrit un profil réaliste avec 10 nombres ; $A_0$ pilote le rayon de nez, donc la SER.
\item Les QUBO de 30 à 33 qubits sont résolus par un recuit réglé sur l'objectif, le meilleur des trois réglages mesurés.
\item On mesure l'hypervolume relatif à un front de référence commun à toutes les méthodes.
\end{itemize}
\end{resultat}

\exercices
\begin{exercice}\label{ex:cst}
(a) Montrer que pour $x\to0$, $y(x)\approx A_0\sqrt x$, et en déduire que le rayon de courbure au nez vaut $A_0^2/2$. (Indice : le cercle $(x-R)^2+y^2=R^2$ donne $y^2\approx2Rx$ près de $x=0$.) (b) Que vaut $y$ au bord de fuite ? Que faudrait-il ajouter pour un bord de fuite épais ?
\end{exercice}
\begin{exercice}\label{ex:echelle-taille}
Combien de temps faudrait-il pour énumérer les $8^{10}$ profils à 1 seconde par évaluation ? Et pour énumérer les $2^{33}$ chaînes d'un QUBO à raison de $10^8$ énergies par seconde ? Laquelle des deux énumérations est hors de portée, et pourquoi l'autre l'est aussi en pratique ?
\end{exercice}

\begin{pistes}
\begin{itemize}
\item Ajouter un troisième niveau de difficulté : 20 coefficients (CST de degré 9) et 60 bits.
\item Remplacer l'optique physique par la méthode des moments pour les seuls profils du front final (vérification haute fidélité).
\end{itemize}
\end{pistes}

\chapter{(à venir)}\label{ch:resultats-echelle}

\chapter{Ce que l'IA peut apporter d'autre}\label{ch:ia-autres}
\iterbanner{V3}

Ce cours a testé une idée précise : apprendre le QUBO. L'apprentissage automatique peut intervenir ailleurs dans la chaîne, et parfois là où le coût se trouve vraiment, c'est-à-dire dans la simulation. Ce chapitre fait le point, à partir d'une recherche bibliographique faite en novembre 2026. Les références sont à vérifier dans les articles avant d'être citées : plusieurs sont des prépublications récentes.

\section{Remplacer le simulateur aérodynamique}
\paragraph{L'idée.} Entraîner un réseau de neurones sur des milliers de calculs CFD pour qu'il prédise $C_L$, $C_D$, voire le champ de pression, en quelques millisecondes. L'optimisation se fait ensuite sur le réseau.
\paragraph{L'état de l'art.} Une revue de 2025 (\emph{Physics of Fluids}) recense les réseaux \og pilotés par les données\fg{} (convolutifs, génératifs) et \og pilotés par la physique\fg{} (PINN, opérateurs neuronaux DeepONet et FNO) utilisés pour l'optimisation de profils et d'aubes. Les opérateurs neuronaux ont été testés comme substituts pour l'optimisation de forme de profils (2023). NeuralFoil (2025) est un outil ouvert, entraîné sur des calculs de type XFOIL, qui donne $C_L$ et $C_D$ quasi instantanément avec une estimation de son incertitude.
\paragraph{Pour nous.} NeuralFoil pourrait servir de simulateur \emph{basse fidélité} gratuit, XFOIL restant la référence. C'est la piste la plus directe pour réduire le coût, et elle ne demande rien au quantique.

\section{Remplacer le calcul de SER}
\paragraph{L'état de l'art.} Des réseaux informés par la physique apprennent les courants induits sur la surface d'une cible métallique 3D, puis en déduisent la SER par l'intégrale de rayonnement, dans toutes les directions (U-PINet, 2025). Les auteurs soulignent que les réseaux sans contrainte physique généralisent mal à des géométries nouvelles. Des réseaux à deux branches ciblent les objets diélectriques 3D (IEEE TAP, 2025).
\paragraph{Pour nous.} C'est pertinent pour le passage au 3D (aile complète), où la méthode des moments devient très coûteuse ; à discuter avec le groupe SER.

\section{Combiner plusieurs fidélités}
\paragraph{L'idée.} On dispose souvent d'un modèle rapide mais approché et d'un modèle lent mais précis. Un modèle multi-fidélité apprend la correction de l'un vers l'autre, et l'optimisation choisit à chaque pas quel niveau évaluer.
\paragraph{L'état de l'art.} Le co-krigeage reste la référence quand les calculs haute fidélité sont rares : une étude de Chalmers (RANS contre MISES) et une étude de 2025 sur une aube de redresseur trouvent le co-krigeage meilleur que les réseaux de neurones multi-fidélité avec quelques dizaines de calculs haute fidélité. Les réseaux deviennent intéressants avec plus de données ou des sorties de grande dimension (champs complets).
\paragraph{Pour nous.} Le projet a déjà une paire naturelle : l'optique physique (rapide) et la méthode des moments (exacte), dont le classement est presque identique (corrélation de rang 0,999). Une boucle multi-fidélité explorerait avec la PO et ne confirmerait que les candidats prometteurs par la MoM. C'est la piste la plus rentable à court terme.

\section{L'optimisation bayésienne en grande dimension}
\paragraph{L'état de l'art.} Deux méthodes de référence : TuRBO (Eriksson et al., 2019), qui remplace un modèle global par plusieurs modèles locaux dans des régions de confiance qui rétrécissent ou redémarrent ; SAASBO (Eriksson et Jankowiak, UAI 2021), un krigeage avec un a priori qui suppose que peu de variables comptent vraiment, et qui les trouve. BAxUS étend TuRBO en augmentant progressivement la dimension du sous-espace. SAASBO coûte cher (inférence en $N^3$) et se limite à de petits budgets, ce qui est justement notre cas.
\paragraph{Pour nous.} C'est le concurrent classique naturel du problème mis à l'échelle : la question \og un modèle global quadratique suffit-il en dimension 20 ?\fg{} a déjà une réponse de la littérature, plutôt négative, et des méthodes locales.

\section{La conception générative et les espaces latents}
\paragraph{L'idée.} Apprendre, sur une base de formes existantes, un espace de petite dimension où chaque point est une forme valide (autoencodeur variationnel, réseau antagoniste génératif, modèle de diffusion), puis optimiser dans cet espace.
\paragraph{L'état de l'art.} Les modèles de diffusion sont la tendance récente pour générer des profils à caractéristiques imposées (portance, traînée) ; un débat porte sur le fait qu'un espace latent n'est qu'une paramétrisation de plus, qui restreint l'espace de conception. Une base de 200\,000 profils (AFBench) permet de comparer ces approches. Côté QUBO, l'autoencodeur binaire de Wilson et al. (2021) donne directement un espace de bits où les contraintes sont respectées par construction.
\paragraph{Pour nous.} C'est la seule piste qui attaque le défaut principal identifié par le projet : les pénalités. Elle demande une base de formes faisables pour entraîner l'autoencodeur, que la grille V2 (702 profils faisables) ou le cache du problème CST pourraient fournir.

\section{L'IA pour le quantique lui-même}
\paragraph{L'état de l'art.} Des réseaux de neurones sur graphes prédisent de bons angles initiaux pour QAOA à partir du graphe du problème (2024) ; d'autres travaux transfèrent les angles d'un problème à un autre et apprennent à prédire quand ce transfert marchera (2025). Le but est de réduire le nombre d'appels au processeur quantique pendant le réglage des angles.
\paragraph{Pour nous.} Nos QUBO successifs se ressemblent d'une itération à l'autre : transférer les angles de QAOA de l'itération précédente est la version simple de cette idée, déjà notée comme piste dans le cours de calcul quantique. Elle réduirait le coût de QAOA sans changer le constat principal : le solveur n'est pas le goulot.

\section{Ce qui compte, par ordre de priorité pour le projet}
\begin{center}
\renewcommand{\arraystretch}{1.25}
\small
\begin{tabularx}{\linewidth}{c>{\raggedright\arraybackslash}X>{\raggedright\arraybackslash}p{3.3cm}>{\raggedright\arraybackslash}p{2.6cm}}
\toprule
& Piste & Gain attendu & Effort \\
\midrule
1 & multi-fidélité PO / MoM & moins de calculs MoM & faible, données déjà là \\
2 & simulateur aéro appris (NeuralFoil) en basse fidélité & moins d'appels XFOIL & faible \\
3 & espace latent binaire pour les contraintes & QUBO sans pénalités, utile au quantique & élevé \\
4 & optimisation bayésienne locale (TuRBO) à l'échelle & concurrent classique honnête & moyen \\
5 & SER apprise en 3D & passage à l'aile complète & élevé, avec le groupe SER \\
6 & angles de QAOA appris ou transférés & coût de QAOA & faible, gain marginal \\
\bottomrule
\end{tabularx}
\end{center}

\begin{resultat}
L'apprentissage automatique a le plus d'impact là où se trouve le coût : la \textbf{simulation} (substituts neuronaux, multi-fidélité). Pour la partie quantique, son apport principal serait de \textbf{supprimer les pénalités} en apprenant un espace de conceptions faisables, ce qui attaque directement le défaut mesuré dans ce projet.
\end{resultat}

\section*{Références de ce chapitre}
\begin{itemize}\small
\item Revue \og deep learning-based aerodynamic shape surrogate models and optimization for airfoils and blade profiles\fg{}, \emph{Physics of Fluids} 37(4), 2025.
\item Deep neural operators as surrogates for shape optimization (airfoils), arXiv~2302.00807.
\item NeuralFoil : airfoil aerodynamics analysis using physics-informed machine learning, arXiv~2503.16323.
\item U-PINet : physics-informed hierarchical learning for 3D RCS prediction, arXiv~2508.03774.
\item Co-kriging contre réseaux multi-fidélité : arXiv~2505.17279 ; étude Chalmers (RANS/MISES), research.chalmers.se, publication 552682.
\item TuRBO (Eriksson et al., NeurIPS 2019) ; SAASBO (Eriksson et Jankowiak, UAI 2021, arXiv~2103.00349).
\item Diffusion pour profils : arXiv~2408.15898 ; AFBench, arXiv~2406.18846.
\item Wilson et al., \emph{Applied Physics Reviews} 8 (2021), arXiv~2105.02396 ; autoencodeurs binaires et FMQA, arXiv~2602.10037.
\item Angles de QAOA par réseaux sur graphes : arXiv~2403.03310 ; transfert d'angles : arXiv~2504.10733.
\end{itemize}

\chapter{Bilan}\label{ch:bilan}
\iterbanner{V3}

\section{La réponse à la question du chapitre~\ref{ch:question}}
La question était : des modèles appris (machine de factorisation, modèle bayésien sur les bits, réseau de neurones) peuvent-ils remplacer le modèle quadratique en $z$, et nous rapprocher de problèmes plus grands et plus « boîtes noires » ? Les deux problèmes donnent des réponses différentes, et c'est la différence qui est instructive.
\begin{center}
\renewcommand{\arraystretch}{1.25}
\small
\begin{tabularx}{\linewidth}{>{\raggedright\arraybackslash}p{3.3cm} >{\raggedright\arraybackslash}X >{\raggedright\arraybackslash}X}
\toprule
& Problème V2 (4 variables, 10 bits) & Problème CST (10 variables, 30 bits) \\
\midrule
modèle en $z$ (actuel) & le meilleur & à égalité avec les autres QUBO \\
machine de factorisation & proche, non significatif & en avance en début de budget, non significatif \\
modèle bayésien sur les bits & moins bon & à égalité ; trouve les extrémités du front \\
réseau seul & échoue (modèle de contrainte trop fidèle) & non testé \\
témoin hybride (non QUBO) & mieux sur S2, moins bien sur S1 & pas significativement mieux \\
contrainte en boîte noire & ne retrouve jamais le front & moins bonne que la contrainte modélisée \\
boucle QUBO contre NSGA-II & meilleure (V2) & \textbf{meilleure}, $p=0{,}008$ à 100 simulations \\
\bottomrule
\end{tabularx}
\end{center}
\begin{itemize}
\item \textbf{À petite dimension}, la connaissance de l'ordre des niveaux (le modèle en $z$) vaut plus que la souplesse : les modèles en bits n'apportent rien.
\item \textbf{À l'échelle}, les trois modèles quadratiques se valent, et tous battent les méthodes sans modèle. Le modèle en $z$ perd son avance parce que ses 66 coefficients demandent autant de données que le budget en fournit.
\item \textbf{Le traitement des contraintes décide plus que le modèle d'objectif}, aux deux échelles : un modèle de contrainte trop sévère exclut le front (réseau seul, V2) ; une contrainte oubliée détourne l'optimisation (première campagne CST) ; une contrainte en boîte noire rend la méthode trop prudente.
\end{itemize}

\section{Pourquoi nos résultats et la littérature divergent, puis se rejoignent}
Les articles sur FMQA et BOCS montrent des gains nets. Sur le problème V2, on n'en voit aucun ; sur le problème CST, on commence à les voir. Les différences expliquent les deux constats :
\begin{center}
\renewcommand{\arraystretch}{1.25}
\small
\begin{tabularx}{\linewidth}{>{\raggedright\arraybackslash}p{2.9cm} >{\raggedright\arraybackslash}X >{\raggedright\arraybackslash}X >{\raggedright\arraybackslash}X}
\toprule
& Littérature FMQA / BOCS & Notre V2 & Notre CST \\
\midrule
variables & binaires par nature (matériaux, graphes) & 4 niveaux ordonnés : l'ordre est une information que le modèle en $z$ exploite & 10 niveaux ordonnés, mais 66 coefficients à apprendre \\
dimension & 30 à 100 bits et plus & 10 bits : tout s'énumère & 30 bits : ni le problème ni le QUBO ne s'énumèrent \\
données par coefficient & souvent moins d'une & environ 5 & environ 1 à 2 \\
front & souvent loin des contraintes & au bord de la contrainte de manœuvre & au bord de plusieurs contraintes \\
mesure & meilleure valeur trouvée & front exact retrouvé (très strict) & hypervolume relatif \\
référence & souvent hasard ou recuit seul & une boucle QUBO déjà bien réglée & la même \\
\bottomrule
\end{tabularx}
\end{center}
La ligne décisive est celle des \textbf{données par coefficient}. Les modèles en bits régularisés sont conçus pour le régime où l'on a moins de données que de coefficients. La V2 n'est pas dans ce régime ; le problème CST s'en approche, et c'est là qu'ils commencent à servir. Un vrai problème de conception (dizaines de paramètres, simulations de plusieurs heures, quelques centaines de calculs au plus) y est entièrement.

Il y a aussi une raison propre à la recherche : le matériel quantique ne sait résoudre que des QUBO. Les travaux cherchent donc d'abord le meilleur modèle \emph{quadratique} ; ils comparent rarement à un modèle libre ou à une méthode classique bien réglée. C'est ce que fait le témoin hybride, et la conclusion de ce cours en dépend.

\section{Ce que cela dit pour le quantique}
\begin{enumerate}
\item \textbf{À l'échelle, la forme QUBO ne coûte rien de mesurable.} Le témoin non quantique ne fait pas significativement mieux que les modèles quadratiques. Écrire le problème comme un QUBO, condition pour le confier à une machine quantique, n'a pas fait perdre de qualité.
\item \textbf{Les QUBO deviennent assez grands pour qu'un solveur matériel ait un sens.} Avec 33 qubits, on ne peut plus énumérer dans la boucle, et l'on résout par recuit simulé. C'est exactement la place d'un recuit quantique. Ce cours \emph{ne montre pas} qu'il ferait mieux : il montre que la question devient légitime.
\item \textbf{Le goulot reste la simulation.} Le choix du modèle change le résultat de quelques points d'hypervolume ; un QUBO mieux résolu, beaucoup moins, puisque le recuit réglé trouve déjà l'optimum de deux tiers des QUBO de la V2 (chapitre~\ref{ch:echelle}), et qu'un candidat proche de l'optimum du modèle est presque aussi bon à simuler. L'apport principal de l'IA est en amont : simulateurs appris, multi-fidélité, espaces latents sans pénalités (chapitre~\ref{ch:ia-autres}).
\end{enumerate}

\section{Comment le présenter}
Pour un public non spécialiste, en trois phrases :
\begin{quote}
« Chaque simulation d'un profil coûte cher ; on apprend donc, à partir des profils déjà simulés, un modèle mathématique qui prédit la traînée et la signature radar, puis on demande au modèle quel profil essayer ensuite. Ce modèle s'écrit sous la forme qu'une machine quantique sait résoudre ; sur un problème de 30 bits, cela fonctionne aussi bien qu'un modèle sans cette contrainte, et mieux qu'un algorithme génétique. Le point délicat n'est pas le calcul quantique, mais les contraintes : une contrainte mal décrite détourne toute l'optimisation. »
\end{quote}
Et une précaution : ne pas présenter les résultats de la première campagne CST, qui donnaient un avantage spectaculaire aux méthodes de la littérature. Cet avantage venait d'un défaut du protocole.

\begin{resultat}
\begin{itemize}
\item Les modèles appris n'apportent rien sur le problème V2 ; sur le problème CST de 30 bits, ils font jeu égal avec le modèle en $z$, et la machine de factorisation est en avance en début de budget.
\item Aux deux échelles, la boucle QUBO bat NSGA-II, et le traitement des contraintes compte plus que le modèle d'objectif.
\item À l'échelle, la forme QUBO ne coûte rien de mesurable par rapport à un modèle libre : c'est l'argument le plus solide de ce cours pour une résolution quantique.
\item La littérature et nos résultats divergent surtout par le nombre de données par coefficient ; notre problème CST se rapproche du régime des articles, et les résultats s'en rapprochent aussi.
\end{itemize}
\end{resultat}

\exercices
\begin{exercice}\label{ex:bilan-feynman}
Expliquer à quelqu'un qui connaît la régression linéaire, mais pas le QUBO, pourquoi le modèle en $z$ est le meilleur sur la V2 et ne l'est plus à 30 bits. On doit pouvoir le faire en cinq phrases, sans le mot « qubit ».
\end{exercice}
\begin{exercice}\label{ex:bilan-defense}
Un ingénieur objecte : « Votre boucle bat NSGA-II, mais NSGA-II n'est pas fait pour 300 simulations. » Que répondre ? Quelle méthode classique faudrait-il ajouter à la comparaison pour que l'argument tienne (voir le chapitre~\ref{ch:ia-autres}) ?
\end{exercice}

\begin{pistes}
\begin{itemize}
\item Ajouter un concurrent classique adapté aux petits budgets en grande dimension (TuRBO ou SAASBO, chapitre~\ref{ch:ia-autres}) au problème CST : c'est la comparaison qui manque pour défendre la boucle.
\item 20 graines sur le problème CST, pour trancher entre la machine de factorisation et le modèle en $z$.
\item Un troisième niveau (CST de degré 9, 60 bits), où le modèle en $z$ aurait 231 coefficients et devrait être distancé par les modèles en bits.
\item Résoudre les QUBO de 33 qubits sur un recuit quantique réel, et comparer au recuit simulé à budget égal.
\end{itemize}
\end{pistes}


\appendix
\chapter{Corrigés des exercices}\label{annexe:corriges}
\iterbanner{V3}

\section*{Chapitre~\ref{ch:question} — La question}
\textbf{Ex.~1.1.} \code{quadz} : un QUBO (après substitution $z=a_0+Ab$), 15 paramètres par grandeur, et il connaît l'ordre des niveaux puisque $z$ est proportionnel à l'indice. \code{fm} : un QUBO par construction, $1+10+10\times4=51$ paramètres, n'a aucune notion d'ordre. \code{bocs} : chaque tirage est un QUBO, $1+10+45=56$ paramètres, pas d'ordre non plus. \code{mlp} : pas un QUBO, 1\,249 paramètres, connaît l'ordre (son entrée est $z$).

\textbf{Ex.~1.2.} La boucle V2 suit exactement le schéma FMQA : apprendre un modèle quadratique en bits, le résoudre comme un QUBO, simuler le candidat, réapprendre. La différence tient à la forme des couplages. Dans le modèle en $z$, ils valent $2(A^\top MA)_{kl}$ : ils sont tous déterminés par les 15 coefficients du polynôme. Par exemple, deux bits $j$ et $l$ d'une même variable sont couplés par $2M_{ii}a_ja_l$, proportionnel à $2^{j+l}$ : la structure binaire est imposée. Dans une machine de factorisation, chaque couplage $\langle\mathbf v_k,\mathbf v_l\rangle$ est libre, à la contrainte de rang près.

\section*{Chapitre~\ref{ch:fm} — Machine de factorisation}
\textbf{Ex.~2.1.} $\hat y(\mathbf b)=1-b_1+0{,}5b_2-2b_3+0{,}5b_1b_2-b_1b_3$. Les huit valeurs ($b_1b_2b_3$) : $000\to1$ ; $100\to0$ ; $010\to1{,}5$ ; $110\to1$ ; $001\to-1$ ; $101\to-3$ ; $011\to-0{,}5$ ; $111\to-2$. Le minimum est $101$ ($-3$), vérifié dans le cours par la formule rapide.

\textbf{Ex.~2.2.} (a) En rang 1, $J_{ij}=v_iv_j$ (nombres réels), donc $J_{12}J_{13}J_{23}=(v_1v_2v_3)^2\ge0$, alors que $(-1)^3=-1$. (b) $\mathbf v_i=\sqrt2(\cos\theta_i,\sin\theta_i)$ avec $\theta_i=0^\circ,120^\circ,240^\circ$ : $\langle\mathbf v_i,\mathbf v_j\rangle=2\cos120^\circ=-1$. (c) Le rang limite les structures de couplages représentables, ici des « frustrations » où trois bits se repoussent deux à deux. Un rang trop petit empêche de représenter certaines physiques ; un rang trop grand ajoute des paramètres. On le règle par validation ; le code prend $k=4$.

\textbf{Ex.~2.3.} $3=110$ et $4=001$ diffèrent par leurs 3 bits ; $0=000$ et $7=111$ aussi ; $6=011$ et $7=111$ par un seul. Le niveau 4 est le seul où le bit $b_2$ est allumé seul. Sa prédiction vaut la constante plus le poids de $b_2$, et ce poids est appris sur les niveaux 5, 6 et 7, où $b_2$ est allumé \emph{avec} d'autres bits et leurs couplages. Sans profil au niveau 4, rien ne contraint la combinaison « $b_2$ seul » : c'est une extrapolation. Le modèle en $z$, lui, impose une parabole en $z$ qui passe nécessairement entre les niveaux 3 et 5 : il interpole.

\section*{Chapitre~\ref{ch:bocs} — Modèle bayésien}
\textbf{Ex.~3.1.} Le gradient de $\beta\|\Phi\boldsymbol\theta-\mathbf y\|^2+\boldsymbol\theta^\top A\boldsymbol\theta$ vaut $2\beta\Phi^\top(\Phi\boldsymbol\theta-\mathbf y)+2A\boldsymbol\theta$. Il s'annule pour $(A+\beta\Phi^\top\Phi)\boldsymbol\theta=\beta\Phi^\top\mathbf y$, soit $\boldsymbol\theta=\beta S\Phi^\top\mathbf y=\mathbf m$. En divisant par $\beta$ : c'est la solution ridge avec $\lambda=\alpha/\beta$ sur tous les coefficients sauf la constante.

\textbf{Ex.~3.2.} (a) $P(\theta<0)=\Phi(-0{,}3/0{,}2)=\Phi(-1{,}5)\approx0{,}07$. (b) $\Phi(-0{,}3)\approx0{,}38$. (c) En début de boucle, peu de données : la covariance $S$ est large, les tirages s'écartent souvent de la moyenne, et des régions jugées moyennes par $\mathbf m$ sont tirées comme bonnes, puis explorées. Quand les données s'accumulent, $S$ rétrécit et les tirages convergent vers $\mathbf m$ : on exploite.

\textbf{Ex.~3.3.} $\Phi$ a $N=30$ lignes, donc un rang au plus 30 : $\Phi^\top\Phi$ a au plus 30 valeurs propres non nulles. Les termes correspondant à $\lambda_k=0$ sont nuls, et chacun des autres est inférieur à 1 : $\gamma\le30$. Pour un modèle à 15 coefficients dont toutes les valeurs propres sont grandes devant $\alpha/\beta$, chaque terme vaut presque 1 : $\gamma\approx15$, tous les coefficients sont déterminés.

\section*{Chapitre~\ref{ch:mlp} — Réseau de neurones}
\textbf{Ex.~4.1.} $dh+h$ (première couche) $+\,h^2+h$ (deuxième) $+\,h+1$ (sortie) $=h^2+(d+3)h+1$. Pour $h=32$ : $1\,249$ ($d=4$) et $1\,441$ ($d=10$). Pour $d=10$, $h^2+13h+1=66$ donne $h\approx4$ (69 paramètres) : un réseau de 4 neurones par couche a autant de paramètres que le modèle quadratique en $z$.

\textbf{Ex.~4.2.} (a) Pour $f=b_0b_1b_2$, seule la chaîne $111$ contribue, avec le signe $(-1)^0=+1$ : le coefficient vaut 1. (b) En sommant avec les signes $(-1)^{3-|\mathbf b|}$ et la correspondance $x=b_0+2b_1+4b_2$, on trouve environ $-0{,}86$ : un terme cubique important, que le QUBO ne peut pas représenter. (c) $z_m$ est affine en bits, donc $z_m^2$ est quadratique : le coefficient cubique est nul (le calcul donne bien 0).

\section*{Chapitre~\ref{ch:precision} — Précision}
\textbf{Ex.~5.1.} $0{,}35\times10{,}4\approx3{,}6$ counts. Les profils voisins du front de S2 sont espacés de 0,1 à 3,9 counts (38,0 ; 40,4 ; 44,3 ; 44,4 ; 47,0 ; 50,9 ; 53,5 ; 55,8 ; 58,7) : l'erreur est du même ordre que l'écart entre voisins. Le modèle situe bien la région du front, mais il ne peut pas départager sûrement deux profils voisins. C'est pourquoi la boucle simule les candidats au lieu de croire le modèle.

\textbf{Ex.~5.2.} La corrélation de rang ne regarde que l'ordre des valeurs, pas leur écart. Un modèle qui prédit $\hat y=y+\sigma$ (la vraie valeur décalée d'un écart-type) classe parfaitement les profils (rang 1) mais se trompe de $\sigma$ partout : erreur relative 1. Pour choisir le meilleur profil, ce modèle serait parfait.

\section*{Chapitre~\ref{ch:contraintes} — Contraintes}
\textbf{Ex.~6.1.} Faisables entre $0{,}3$ et $1{,}1$ : les infaisables reçoivent $1{,}1+0{,}2\times0{,}8=1{,}26$. Si un nouveau faisable vaut $1{,}5$, le maximum et l'étendue changent : les infaisables valent maintenant $1{,}5+0{,}2\times1{,}2=1{,}74$. Les données des infaisables bougent d'une itération à l'autre ; le modèle apprend une cible qui se déplace.

\textbf{Ex.~6.2.} Problème V2 : 14 qubits contre 11, soit 3 qubits d'écart en moins, et $30+3=33$ couplages en moins (chaque bit d'écart avec les 10 bits de conception, et entre eux). Problème mis à l'échelle : 33 contre 30 qubits, et $3\times30+3=93$ couplages en moins.

\section*{Chapitre~\ref{ch:echelle} — Problème mis à l'échelle}
\textbf{Ex.~8.1.} (a) Près de $x=0$, $C(x)\approx\sqrt x$ et $S(0)=A_0$ (seul le premier polynôme de Bernstein vaut 1 en 0), donc $y\approx A_0\sqrt x$ et $y^2\approx A_0^2x$. Le cercle de rayon $R$ tangent en 0 donne $y^2\approx2Rx$ : $R=A_0^2/2$. (b) $C(1)=0$, donc $y(1)=0$ : le bord de fuite est fermé. Pour un bord de fuite épais, on ajoute à chaque face un terme $x\,\Delta y_{\rm BF}/2$.

\textbf{Ex.~8.2.} $8^{10}\approx1{,}07\cdot10^9$ secondes, soit environ 34 ans de simulations : hors de portée. $2^{33}\approx8{,}6\cdot10^9$ énergies à $10^8$ par seconde : environ 90 secondes, ce qui est faisable pour un QUBO. Mais la boucle en résout environ 300 par exécution (5 par itération), soit plus de 7 heures, et il faudrait stocker ou trier des milliards de valeurs. Ce qui rend le problème « grand », c'est la simulation de chaque profil ; le QUBO reste abordable bien plus longtemps. C'est encore le message du cours de calcul quantique : le coût est dans la simulation.

\section*{Chapitre~\ref{ch:resultats-v2} — Résultats V2}
\textbf{Ex.~7.1.} En médiane, $95-75=20$ simulations de plus (la médiane des écarts graine par graine vaut 22,5). Pour une graine donnée, les deux méthodes partent des \emph{mêmes} 10 profils initiaux : une graine « difficile » l'est pour les deux. Comparer graine par graine retire cette variabilité commune et ne garde que l'effet de la méthode. Le test apparié est donc plus sensible que la comparaison de deux médianes calculées sur des tirages qui varient beaucoup.

\textbf{Ex.~7.2.} $75\times(1{,}2-0{,}89)^2=75\times0{,}31^2\approx7{,}2$, soit environ sept fois l'étendue de l'objectif normalisé : aucun gain d'objectif ne peut compenser, et le profil n'est jamais proposé. Pour qu'il le soit, il faudrait que la pénalité reste de l'ordre d'un dixième de l'étendue, soit $1{,}2-\widehat{C_{L,\max}}\lesssim\sqrt{0{,}1/75}\approx0{,}04$ : une erreur de prédiction de moins de 0,04 sur une grandeur irrégulière. Autre solution : relâcher la contrainte là où le modèle est incertain (probabilité de faisabilité), ou utiliser un modèle plus indulgent pour la contrainte seule.

\section*{Chapitre~\ref{ch:resultats-echelle} — Résultats à l'échelle}
\textbf{Ex.~9.1.} (a) Sous l'hypothèse nulle, chaque différence a autant de chances d'être positive que négative, et les $2^8=256$ configurations de signes sont équiprobables. La plus extrême est celle où les 8 différences ont le même signe (statistique de Wilcoxon nulle) ; en bilatéral, on compte les deux sens : $p=2/256\approx0{,}0078$. Aucune configuration ne peut donner moins avec 8 graines. (b) Test des signes : $P(X\ge6)=\big(\binom86+\binom87+\binom88\big)/256=(28+8+1)/256\approx0{,}145$, donc $p\approx0{,}29$ en bilatéral. (c) Avec 20 graines et 15 gagnées : $P(X\ge15)=(15\,504+4\,845+1\,140+190+20+1)/2^{20}\approx0{,}021$, soit $p\approx0{,}041<0{,}05$. Avec 16 graines et 12 gagnées, $p\approx0{,}077$ : il faut une vingtaine de graines.

\textbf{Ex.~9.2.} (a) $300\times(1-0{,}59-0{,}14)=81$ simulations admissibles par graine. (b) 36 profils faisables pour 66 coefficients : le système est sous-déterminé. Avec $\lambda=10^{-6}$, le modèle en $z$ prend la solution de norme minimale (cours de calcul quantique, chapitre~8), qui passe par les données mais dit n'importe quoi ailleurs ; on est avant le pic « juste déterminé » du chapitre~\ref{ch:precision}, où la corrélation de rang sur $C_D$ reste vers 0,37. (c) La régularisation (pénalité $L_2$, a priori bayésien) ramène vers zéro les coefficients que les données ne déterminent pas. Dans une région inexplorée, les prédictions restent donc modérées au lieu d'extrapoler vers de très bonnes valeurs : le modèle propose moins de profils extrêmes, et moins de géométries invalides.

\section*{Chapitre~\ref{ch:bilan} — Bilan}
\textbf{Ex.~11.1.} Une réponse possible : « On approche la traînée par un polynôme de degré 2 des paramètres de forme, et l'on choisit le prochain profil en minimisant ce polynôme. Avec 4 paramètres, le polynôme a 15 coefficients, et après 75 simulations on en a cinq fois plus que de coefficients : il est bien déterminé, et le fait de savoir que les niveaux d'un paramètre se suivent est une information gratuite. Avec 10 paramètres, il a 66 coefficients, et l'on n'a guère plus de profils faisables que de coefficients : le polynôme est mal déterminé et peut extrapoler n'importe où. Les autres modèles ajoutent une régularisation, qui les rend prudents quand les données manquent. Ils rattrapent donc leur retard quand la dimension augmente. »

\textbf{Ex.~11.2.} L'objection est juste : avec une population de 16 et 300 simulations, NSGA-II fait moins de 20 générations, ce qui est peu pour un algorithme évolutionnaire. Battre NSGA-II montre que le modèle apporte de l'information, pas que la boucle bat les meilleures méthodes classiques. Il faut ajouter un concurrent conçu pour les petits budgets : une optimisation bayésienne, adaptée à la grande dimension (TuRBO, SAASBO). Sur la V2, le lot L1 a déjà montré qu'une optimisation bayésienne fait jeu égal avec la boucle. L'argument honnête est donc : « la boucle QUBO est au niveau des bonnes méthodes à modèle, et sa forme la rend compatible avec une machine quantique ».


\end{document}
